%=============================================================
%  TIERED BIBLIOGRAPHY  (v2 — QM only, expanded)
%  Quantum Mechanics & Quantum Chemistry Programme
%  One section per learning tier (HS → PhD) + Full Catalogue
%=============================================================
\documentclass[11pt,a4paper]{article}

\usepackage[T1]{fontenc}
\usepackage[utf8]{inputenc}
\usepackage{fix-cm}
\usepackage[margin=2.5cm]{geometry}
\usepackage{amsmath,amssymb}
\usepackage{booktabs}
\usepackage{array}
\usepackage{longtable}
\usepackage{tabularx}
\usepackage{enumitem}
\usepackage{sectsty}
\usepackage{fancyhdr}
\usepackage[dvipsnames,svgnames,table]{xcolor}
\usepackage{hyperref}
\usepackage{parskip}
\usepackage{tcolorbox}
\tcbuselibrary{skins,breakable}

%--- Colours -------------------------------------------------
\definecolor{darkblue}{RGB}{10,40,90}
\definecolor{midblue}{RGB}{30,80,160}
\definecolor{lightblue}{RGB}{220,233,255}
\definecolor{accentgold}{RGB}{180,140,20}
\definecolor{lightgold}{RGB}{255,248,220}
\definecolor{darkgray}{RGB}{60,60,60}
\definecolor{t1green}{RGB}{30,120,50}
\definecolor{t1light}{RGB}{212,242,220}
\definecolor{t2blue}{RGB}{25,90,170}
\definecolor{t2light}{RGB}{215,232,255}
\definecolor{t3purple}{RGB}{110,35,160}
\definecolor{t3light}{RGB}{238,218,255}
\definecolor{t4orange}{RGB}{165,70,10}
\definecolor{t4light}{RGB}{255,232,208}
\definecolor{t5red}{RGB}{140,15,15}
\definecolor{t5light}{RGB}{255,215,215}
\definecolor{catgray}{RGB}{80,80,80}
\definecolor{catlight}{RGB}{240,240,240}

%--- Section fonts -------------------------------------------
\sectionfont{\Large\bfseries\color{darkblue}\sectionrule{0ex}{0pt}{-1ex}{1.5pt}}
\subsectionfont{\large\bfseries\color{midblue}}
\subsubsectionfont{\normalsize\bfseries\color{darkgray}}

%--- Hyperref ------------------------------------------------
\hypersetup{colorlinks=true,linkcolor=midblue,urlcolor=midblue,
  pdftitle={QM Programme --- Tiered Bibliography v2}}

%--- Header / Footer -----------------------------------------
\pagestyle{fancy}\fancyhf{}
\fancyhead[L]{\small\color{darkgray}\textit{QM \& QC Programme --- Tiered Bibliography}}
\fancyhead[R]{\small\color{darkgray}\thepage}
\fancyfoot[C]{\small\color{darkgray}\textit{Companion to the Pedagogical Support Material}}
\renewcommand{\headrulewidth}{0.4pt}

%--- tcolorbox styles ----------------------------------------
\tcbset{
  tierheader/.style={enhanced,colback=#1,colframe=#1!80!black,
    boxrule=0pt,arc=0pt,sharp corners,
    left=8pt,right=8pt,top=6pt,bottom=6pt,
    fontupper=\bfseries\large\color{white}},
  entrybox/.style={enhanced,breakable,colback=#1,colframe=#1!60!black,
    boxrule=0.5pt,sharp corners,left=8pt,right=8pt,top=5pt,bottom=5pt},
  annotbox/.style={enhanced,breakable,colback=lightgold,colframe=accentgold,
    boxrule=0.6pt,sharp corners,left=5pt,right=5pt,top=3pt,bottom=3pt,
    fontupper=\small\itshape\color{darkgray}},
  keybox/.style={enhanced,colback=lightblue,colframe=midblue,
    boxrule=0.5pt,rounded corners,left=5pt,right=5pt,top=4pt,bottom=4pt,
    fontupper=\small},
  catbox/.style={enhanced,breakable,colback=catlight,colframe=catgray,
    boxrule=0.6pt,sharp corners,left=6pt,right=6pt,top=5pt,bottom=5pt},
}

\newcommand{\tierbadge}[2]{\colorbox{#1}{\small\bfseries\color{white}\strut\ #2\ }}
\newcommand{\essential}{$\bigstar\bigstar\bigstar$}
\newcommand{\recommended}{$\bigstar\bigstar$}
\newcommand{\useful}{$\bigstar$}
\newenvironment{bibentry}[1]{\begin{tcolorbox}[entrybox=#1,top=4pt,bottom=4pt]}{\end{tcolorbox}\smallskip}

%=============================================================
\begin{document}
%=============================================================

\begin{titlepage}
\pagecolor{darkblue}\color{white}
\vspace*{2.5cm}
\begin{center}
  {\fontsize{28}{34}\selectfont\bfseries Tiered Bibliography}\\[0.4cm]
  {\fontsize{17}{22}\selectfont Quantum Mechanics \& Quantum Chemistry Programme}\\[0.5cm]
  {\color{accentgold}\rule{10cm}{2pt}}\\[0.5cm]
  {\large\itshape Quantum Mechanics --- Reading Lists by Learning Tier}\\[1.8cm]
  \begin{tabular}{@{}cl@{}}
    \colorbox{t1green}{\color{white}\bfseries Tier~1} & High School --- Divulgation \& First Introductions\\[3pt]
    \colorbox{t2blue}{\color{white}\bfseries Tier~2}  & Beginning Undergraduate --- Standard Textbooks\\[3pt]
    \colorbox{t3purple}{\color{white}\bfseries Tier~3} & Advanced Undergraduate --- Rigorous Texts \& Historical Papers\\[3pt]
    \colorbox{t4orange}{\color{white}\bfseries Tier~4} & Master's --- Advanced Textbooks, Monographs \& Research Papers\\[3pt]
    \colorbox{t5red}{\color{white}\bfseries Tier~5}   & PhD --- Research Monographs, Specialised Topics \& Primary Literature\\[6pt]
    \colorbox{catgray}{\color{white}\bfseries Cat.~A--G} & Full Alphabetical Catalogue (all tiers)\\
  \end{tabular}\\[1.8cm]
  {\color{accentgold}\rule{10cm}{1pt}}\\[0.4cm]
  {\small\itshape\essential\ Essential \quad \recommended\ Strongly recommended \quad \useful\ Useful supplement\\[3pt]
  \textbf{[N]} Popular/narrative \quad \textbf{[T]} Textbook \quad
  \textbf{[P]} Primary paper \quad \textbf{[M]} Monograph \quad \textbf{[R]} Review}
\end{center}
\vfill
\end{titlepage}
\nopagecolor

\tableofcontents
\newpage

%=============================================================
\section*{How to Use This Bibliography}
\addcontentsline{toc}{section}{How to Use This Bibliography}
%=============================================================

\begin{tcolorbox}[keybox]
This bibliography accompanies the \textbf{Pedagogical Support Material} and the
\textbf{Course Syllabus} of the QM \& QC Programme. It is purely quantum-mechanical
in scope, covering quantum mechanics, quantum chemistry, and their mathematical
foundations. Reading lists are stratified by the same five tiers used for content
delivery and assessment. The tiers are cumulative: Tier~$k$ reading is assumed as
background for Tier~$k+1$.
\end{tcolorbox}

\begin{center}
\begin{tabular}{clp{3.2cm}p{6cm}}
\toprule
\textbf{Tier} & \textbf{Label} & \textbf{PSM track} & \textbf{Primary reading mode}\\
\midrule
1 & HS       & Undergrad & Popular books cover-to-cover before the course\\
2 & BegUG    & Undergrad & One textbook chapter per lecture; do exercises\\
3 & AdvUG    & Undergrad & Textbook + primary paper abstracts and introductions\\
4 & MSc      & Graduate  & Primary papers + advanced textbooks as main reading\\
5 & PhD      & Grad + Research & Monographs + full primary literature + arXiv\\
\bottomrule
\end{tabular}
\end{center}

\newpage

%#############################################################
\section{Tier 1 --- Divulgation, Popular Science \& First Introductions}
\label{sec:t1}
%#############################################################

\begin{tcolorbox}[tierheader=t1green]
Tier 1 --- High School (HS): Popular Science \& Divulgation Books
\end{tcolorbox}

\medskip
\begin{tcolorbox}[keybox]
\textbf{Goal:} Build vivid physical intuition and historical awareness.
No mathematical prerequisites beyond arithmetic and basic algebra.
Read these books \emph{before} any formal course to establish the mental pictures
that make the formalism meaningful. They answer ``What is quantum mechanics and
why does it matter?'' not ``How does it work mathematically?''
\end{tcolorbox}

%------------------------------------------------------------
\subsection{Divulgation \& Popular Science}
%------------------------------------------------------------

\begin{bibentry}{t1light}
\textbf{Feynman, R.~P.}\\
\textit{QED: The Strange Theory of Light and Matter}\\
Princeton University Press, 1985.\\[3pt]
\essential\ [\textbf{N}]\\[3pt]
\begin{tcolorbox}[annotbox]
The finest popular account of quantum amplitudes and the path-integral picture.
Feynman explains probability amplitudes as ``arrows'' that rotate and combine ---
exactly the pre-formal picture for Module~I.2 Lectures~L01--L02. Essential
pre-reading for anyone entering the programme.
\end{tcolorbox}
\end{bibentry}

\begin{bibentry}{t1light}
\textbf{Feynman, R.~P., Leighton, R.~B.\ \& Sands, M.}\\
\textit{The Feynman Lectures on Physics, Vol.~III: Quantum Mechanics}\\
Addison-Wesley, 1965. Freely at \texttt{feynmanlectures.caltech.edu}\\[3pt]
\essential\ [\textbf{N/T}]\\[3pt]
\begin{tcolorbox}[annotbox]
Chapters~1--4 (interference, amplitude, identical particles) are accessible at Tier~1
with no calculus. These chapters provide the pre-formal motivation for Module~I.2
Lectures L01--L04 and L09. A free, authoritative resource.
\end{tcolorbox}
\end{bibentry}

\begin{bibentry}{t1light}
\textbf{Cox, B.\ \& Forshaw, J.}\\
\textit{The Quantum Universe: Everything That Can Happen Does Happen}\\
Da Capo Press, 2011.\\[3pt]
\recommended\ [\textbf{N}]\\[3pt]
\begin{tcolorbox}[annotbox]
A carefully accurate popular account covering superposition, entanglement, and the
exclusion principle. Recommended as complementary Tier~1 reading alongside Feynman's
\textit{QED}.
\end{tcolorbox}
\end{bibentry}

\begin{bibentry}{t1light}
\textbf{Kumar, M.}\\
\textit{Quantum: Einstein, Bohr, and the Great Debate about the Nature of Reality}\\
Norton, 2008.\\[3pt]
\recommended\ [\textbf{N}]\\[3pt]
\begin{tcolorbox}[annotbox]
Historical narrative of the foundational debates: Copenhagen interpretation,
EPR paradox, Schr\"odinger's cat. Essential context for understanding why
measurement theory (Module~I.2 Lecture~L07, Module~I.3 Lecture~L10) was so
controversial.
\end{tcolorbox}
\end{bibentry}

\begin{bibentry}{t1light}
\textbf{Gamow, G.}\\
\textit{Thirty Years That Shook Physics: The Story of Quantum Theory}\\
Dover, 1966.\\[3pt]
\recommended\ [\textbf{N}]\\[3pt]
\begin{tcolorbox}[annotbox]
A charming historical account by a participant. Covers Planck's quantum,
Bohr's atom, de Broglie, Schr\"odinger, Heisenberg, and Dirac with light
but accurate physics. Excellent for historical context before Module~I.1.
\end{tcolorbox}
\end{bibentry}

\begin{bibentry}{t1light}
\textbf{Penrose, R.}\\
\textit{The Emperor's New Mind}, Chapters~5--6\\
Oxford University Press, 1989.\\[3pt]
\useful\ [\textbf{N}]\\[3pt]
\begin{tcolorbox}[annotbox]
Philosophical context for quantum mechanics, wave functions, and measurement.
Chapters~5--6 are accessible at Tier~1. Provides useful philosophical grounding
before Module~I.3 Lectures L05--L07.
\end{tcolorbox}
\end{bibentry}

\begin{bibentry}{t1light}
\textbf{Al-Khalili, J.}\\
\textit{Quantum: A Guide for the Perplexed}\\
Weidenfeld \& Nicolson, 2003.\\[3pt]
\recommended\ [\textbf{N}]\\[3pt]
\begin{tcolorbox}[annotbox]
Beautifully illustrated popular account. Covers wave--particle duality, uncertainty,
entanglement, and quantum technology with minimal formalism. Effective visual
introduction for Tier~1 students.
\end{tcolorbox}
\end{bibentry}

\begin{bibentry}{t1light}
\textbf{Rae, A.~I.~M.}\\
\textit{Quantum Physics: Illusion or Reality?} (2nd ed.)\\
Cambridge University Press, 2004.\\[3pt]
\recommended\ [\textbf{N}]\\[3pt]
\begin{tcolorbox}[annotbox]
Focused on the conceptual and philosophical issues: the measurement problem,
hidden variables, many worlds, and Bell's theorem. Excellent bridge between
popular science and the philosophical content of Module~I.2 Lecture~L07.
\end{tcolorbox}
\end{bibentry}

\begin{bibentry}{t1light}
\textbf{Hey, T.\ \& Walters, P.}\\
\textit{The New Quantum Universe} (revised ed.)\\
Cambridge University Press, 2003.\\[3pt]
\useful\ [\textbf{N}]\\[3pt]
\begin{tcolorbox}[annotbox]
Visually rich survey of quantum mechanics and its applications (lasers, transistors,
MRI, cryptography). Good for motivating why quantum mechanics matters before
the formalism of Series~I.
\end{tcolorbox}
\end{bibentry}

\begin{bibentry}{t1light}
\textbf{Rovelli, C.}\\
\textit{Helgoland: Making Sense of the Quantum Revolution}\\
Riverhead Books, 2021.\\[3pt]
\useful\ [\textbf{N}]\\[3pt]
\begin{tcolorbox}[annotbox]
A philosophical-historical account advocating the relational interpretation.
Stimulating for students curious about interpretations of quantum mechanics;
provides a counterpoint to Copenhagen before Module~I.3 Lecture~L10.
\end{tcolorbox}
\end{bibentry}

\newpage
%#############################################################
\section{Tier 2 --- Standard Undergraduate Textbooks}
\label{sec:t2}
%#############################################################

\begin{tcolorbox}[tierheader=t2blue]
Tier 2 --- Beginning Undergraduate (BegUG): Standard Textbooks
\end{tcolorbox}

\medskip
\begin{tcolorbox}[keybox]
\textbf{Goal:} Develop working fluency in quantum mechanics notation and methods.
Assumes first-year calculus and linear algebra. Work through one chapter per
lecture, solving end-of-chapter exercises alongside SPSU problem sets from the PSM.
\end{tcolorbox}

%------------------------------------------------------------
\subsection{Core Quantum Mechanics Textbooks}
%------------------------------------------------------------

\begin{bibentry}{t2light}
\textbf{Griffiths, D.~J.\ \& Schroeter, D.~F.}\\
\textit{Introduction to Quantum Mechanics} (3rd ed.)\\
Cambridge University Press, 2018.\\[3pt]
\essential\ [\textbf{T}]\\[3pt]
\begin{tcolorbox}[annotbox]
The standard first course globally. Develops wave mechanics before Dirac notation.
Chapters~1--3 (wavefunction, Schr\"odinger equation, formalism) parallel Module~I.3
Lectures~L02--L06; Chapter~3 introduces Hilbert space and bra-ket, paralleling
Module~I.2 Lectures~L01--L06. Superb problem sets at the SPSU level.
\end{tcolorbox}
\end{bibentry}

\begin{bibentry}{t2light}
\textbf{Shankar, R.}\\
\textit{Principles of Quantum Mechanics} (2nd ed.)\\
Springer, 1994.\\[3pt]
\essential\ [\textbf{T}]\\[3pt]
\begin{tcolorbox}[annotbox]
Begins with a self-contained linear algebra review (Chapter~1) --- ideal for
bridging Tier~1 to Tier~2. More algebraic than Griffiths; a natural bridge to
Sakurai. Chapter~1 pairs with Module~I.1; Chapters~4--9 support Module~I.2
Lectures~L03--L08.
\end{tcolorbox}
\end{bibentry}

\begin{bibentry}{t2light}
\textbf{Townsend, J.~S.}\\
\textit{A Modern Approach to Quantum Mechanics} (2nd ed.)\\
University Science Books, 2012.\\[3pt]
\recommended\ [\textbf{T}]\\[3pt]
\begin{tcolorbox}[annotbox]
Starts from spin-$\tfrac{1}{2}$ and Dirac notation from chapter one (Sakurai-style).
Excellent for building bra-ket fluency early. Module~I.2 uses this as a parallel
text for Lectures~L01--L04 and L10.
\end{tcolorbox}
\end{bibentry}

\begin{bibentry}{t2light}
\textbf{Binney, J.\ \& Skinner, D.}\\
\textit{The Physics of Quantum Mechanics}\\
Oxford University Press, 2014. Freely at \texttt{www-thphys.physics.ox.ac.uk}\\[3pt]
\recommended\ [\textbf{T}]\\[3pt]
\begin{tcolorbox}[annotbox]
Uses bra-ket from Chapter~1. Excellent treatment of measurement and operators.
Good parallel text for Module~I.2 Lectures~L03--L07. Free access is a strong advantage.
\end{tcolorbox}
\end{bibentry}

\begin{bibentry}{t2light}
\textbf{French, A.~P.\ \& Taylor, E.~F.}\\
\textit{An Introduction to Quantum Physics}\\
Norton, 1978.\\[3pt]
\recommended\ [\textbf{T}]\\[3pt]
\begin{tcolorbox}[annotbox]
An MIT undergraduate text emphasising conceptual clarity and connection to experiment.
Chapters on wave packets, uncertainty, and the hydrogen atom are particularly well
motivated. Good supplementary reading for Module~I.3 Lectures~L02, L03, L11.
\end{tcolorbox}
\end{bibentry}

\begin{bibentry}{t2light}
\textbf{McMahon, D.}\\
\textit{Quantum Mechanics Demystified} (2nd ed.)\\
McGraw-Hill, 2005.\\[3pt]
\useful\ [\textbf{T}]\\[3pt]
\begin{tcolorbox}[annotbox]
A problem-solving companion with worked examples at every step. Not a substitute for
a proper textbook but useful alongside Griffiths for students who need more
step-by-step worked solutions. Supports SPSU practice.
\end{tcolorbox}
\end{bibentry}

\begin{bibentry}{t2light}
\textbf{Phillips, A.~C.}\\
\textit{Introduction to Quantum Mechanics}\\
Wiley, 2003.\\[3pt]
\useful\ [\textbf{T}]\\[3pt]
\begin{tcolorbox}[annotbox]
A concise undergraduate text well suited to a one-semester course. Good on the
physical motivation for wave mechanics and the connection between classical and
quantum descriptions. Supports Module~I.3 Lectures~L05--L07.
\end{tcolorbox}
\end{bibentry}

%------------------------------------------------------------
\subsection{Mathematics Support}
%------------------------------------------------------------

\begin{bibentry}{t2light}
\textbf{Strang, G.}\\
\textit{Introduction to Linear Algebra} (5th ed.)\\
Wellesley--Cambridge Press, 2016.\\[3pt]
\essential\ [\textbf{T}]\\[3pt]
\begin{tcolorbox}[annotbox]
The prerequisite linear algebra text. Chapters~5--6 (eigenvalues, orthogonality)
are directly needed for quantum mechanics. Work through before or alongside Module~I.1.
\end{tcolorbox}
\end{bibentry}

\begin{bibentry}{t2light}
\textbf{Boas, M.~L.}\\
\textit{Mathematical Methods in the Physical Sciences} (3rd ed.)\\
Wiley, 2006.\\[3pt]
\recommended\ [\textbf{T}]\\[3pt]
\begin{tcolorbox}[annotbox]
The standard mathematical methods text for physics undergraduates. Chapter~3
(linear algebra), Chapter~6 (complex analysis), and Chapter~7 (Fourier series)
are directly used in Module~I.2 Lectures~L02, L05, L06.
\end{tcolorbox}
\end{bibentry}

\begin{bibentry}{t2light}
\textbf{Byron, F.~W.\ \& Fuller, R.~W.}\\
\textit{Mathematics of Classical and Quantum Physics}\\
Dover, 1992.\\[3pt]
\recommended\ [\textbf{T}]\\[3pt]
\begin{tcolorbox}[annotbox]
Covers Hilbert spaces, operators, distributions, and Green's functions at a level
accessible to advanced Tier~2 students. Useful mathematical supplement for
Module~I.2 Lectures~L06 and L12.
\end{tcolorbox}
\end{bibentry}

%------------------------------------------------------------
\subsection{Quantum Chemistry Introduction}
%------------------------------------------------------------

\begin{bibentry}{t2light}
\textbf{Jensen, F.}\\
\textit{Introduction to Computational Chemistry} (3rd ed.)\\
Wiley, 2017.\\[3pt]
\recommended\ [\textbf{T}]\\[3pt]
\begin{tcolorbox}[annotbox]
Accessible entry to computational quantum chemistry. Chapters~1--3 provide the Tier~2
background for Series~III and should be read alongside Module~III.1.
\end{tcolorbox}
\end{bibentry}

\begin{bibentry}{t2light}
\textbf{Levine, I.~N.}\\
\textit{Quantum Chemistry} (7th ed.)\\
Pearson, 2014.\\[3pt]
\recommended\ [\textbf{T}]\\[3pt]
\begin{tcolorbox}[annotbox]
The standard undergraduate quantum chemistry textbook. Clear development of
the Schr\"odinger equation, variational method, and perturbation theory before
moving to molecular orbital theory. Supports Module~III.1 and early Module~III.2.
\end{tcolorbox}
\end{bibentry}

\begin{bibentry}{t2light}
\textbf{Atkins, P.~W.\ \& Friedman, R.}\\
\textit{Molecular Quantum Mechanics} (5th ed.)\\
Oxford University Press, 2011.\\[3pt]
\recommended\ [\textbf{T}]\\[3pt]
\begin{tcolorbox}[annotbox]
Bridges undergraduate physical chemistry and quantum mechanics. Excellent treatment
of angular momentum, hydrogen atom, and molecular orbital theory. Supports
Module~II.3 and Module~III.2 Lectures~L02--L08.
\end{tcolorbox}
\end{bibentry}

\newpage
%#############################################################
\section{Tier 3 --- Rigorous Undergraduate, Bridge Texts \& Historical Papers}
\label{sec:t3}
%#############################################################

\begin{tcolorbox}[tierheader=t3purple]
Tier 3 --- Advanced Undergraduate (AdvUG): Rigorous Texts, Bridge Books \& First Primary Papers
\end{tcolorbox}

\medskip
\begin{tcolorbox}[keybox]
\textbf{Goal:} Understand internal consistency, proofs, and the geometric structure
of the formalism. Begin reading primary papers at the abstract and introduction level.
Assumes mathematical analysis and linear algebra at a proof-writing level.
\end{tcolorbox}

%------------------------------------------------------------
\subsection{Core Texts}
%------------------------------------------------------------

\begin{bibentry}{t3light}
\textbf{Sakurai, J.~J.\ \& Napolitano, J.}\\
\textit{Modern Quantum Mechanics} (3rd ed.)\\
Cambridge University Press, 2020.\\[3pt]
\essential\ [\textbf{T}]\\[3pt]
\begin{tcolorbox}[annotbox]
The primary text for Module~I.2. Begins with bra-ket notation and the Stern--Gerlach
experiment. Chapters~1--3 (state vectors, dynamics, angular momentum) are the core of
Series~I. The Sakurai ``operational'' track is one of the two tracks in the Module~I.2
dual pedagogy.
\end{tcolorbox}
\end{bibentry}

\begin{bibentry}{t3light}
\textbf{Dirac, P.~A.~M.}\\
\textit{The Principles of Quantum Mechanics} (4th ed.)\\
Oxford University Press, 1958 (reprinted 1981).\\[3pt]
\essential\ [\textbf{T/P}]\\[3pt]
\begin{tcolorbox}[annotbox]
The original source for bra-ket notation. Chapter~I (``The Principle of
Superposition'') and Chapter~II (``Dynamical Variables and Observables'') must be
read alongside Module~I.2. This is the ``Dirac abstract'' track of the dual pedagogy.
A primary source that functions as a textbook.
\end{tcolorbox}
\end{bibentry}

\begin{bibentry}{t3light}
\textbf{Messiah, A.}\\
\textit{Quantum Mechanics} (2 vols., translated from French)\\
North-Holland, 1961--1962 (Dover reprint 2014).\\[3pt]
\essential\ [\textbf{T}]\\[3pt]
\begin{tcolorbox}[annotbox]
One of the great classical treatises. Volume~I covers wave mechanics, the Schr\"odinger
equation, and the hydrogen atom with extraordinary clarity and rigour. Volume~II
covers angular momentum, perturbation theory, scattering, and identical particles.
The appendices on Hilbert space and group theory are among the best introductions
available. Supports all of Series~I and Module~II.3.
\end{tcolorbox}
\end{bibentry}

\begin{bibentry}{t3light}
\textbf{Gasiorowicz, S.}\\
\textit{Quantum Physics} (3rd ed.)\\
Wiley, 2003.\\[3pt]
\recommended\ [\textbf{T}]\\[3pt]
\begin{tcolorbox}[annotbox]
A rigorous and physically well-motivated undergraduate text. Excellent on the
connection between classical and quantum physics, with a thorough treatment of
perturbation theory and scattering. Supports Module~I.3 Lectures~L09--L11.
\end{tcolorbox}
\end{bibentry}

\begin{bibentry}{t3light}
\textbf{Isham, C.~J.}\\
\textit{Lectures on Quantum Theory: Mathematical and Structural Foundations}\\
Imperial College Press, 1995.\\[3pt]
\recommended\ [\textbf{T}]\\[3pt]
\begin{tcolorbox}[annotbox]
Bridges undergraduate and graduate levels. Rigorous treatment of Hilbert spaces,
operators, and measurement without full functional analysis. Ideal for Tier~3
students preparing for Tier~4; directly supports Module~I.2 Lectures~L11--L12.
\end{tcolorbox}
\end{bibentry}

\begin{bibentry}{t3light}
\textbf{Hannabuss, K.}\\
\textit{An Introduction to Quantum Theory}\\
Oxford University Press, 1997.\\[3pt]
\recommended\ [\textbf{T}]\\[3pt]
\begin{tcolorbox}[annotbox]
Mathematically precise undergraduate text. Good on spectral theory, uncertainty,
and density operators without requiring measure theory. Supports Module~I.2
Lectures~L04, L07, and L11.
\end{tcolorbox}
\end{bibentry}

\begin{bibentry}{t3light}
\textbf{Auletta, G., Fortunato, M.\ \& Parisi, G.}\\
\textit{Quantum Mechanics}\\
Cambridge University Press, 2009.\\[3pt]
\recommended\ [\textbf{T}]\\[3pt]
\begin{tcolorbox}[annotbox]
A comprehensive modern undergraduate-to-graduate text covering both formalism and
interpretation. Notable for its treatment of quantum information, entanglement, and
decoherence at an accessible level. Supports Module~I.2 Lectures~L09--L11 and
Module~I.3 Lecture~L10.
\end{tcolorbox}
\end{bibentry}

\begin{bibentry}{t3light}
\textbf{Peres, A.}\\
\textit{Quantum Theory: Concepts and Methods}\\
Kluwer, 1993 (Springer reprint 2002).\\[3pt]
\recommended\ [\textbf{T}]\\[3pt]
\begin{tcolorbox}[annotbox]
A masterful rigorous treatment with a strong emphasis on the conceptual
foundations, Bell inequalities, and the measurement problem. The treatment of
EPR and entanglement is particularly excellent. Supports Module~I.2 Lectures~L07,
L09, and L11 from a foundational perspective.
\end{tcolorbox}
\end{bibentry}

%------------------------------------------------------------
\subsection{Historical Primary Papers --- First Encounters}
%------------------------------------------------------------

\begin{bibentry}{t3light}
\textbf{Heisenberg, W.}\\
\textit{\"Uber quantentheoretische Umdeutung kinematischer und mechanischer
Beziehungen}\\
Z.~Phys.~\textbf{33}, 879--893 (1925). English translation in van der Waerden (ed.),
\textit{Sources of Quantum Mechanics}, Dover, 1967.\\[3pt]
\recommended\ [\textbf{P}]\\[3pt]
\begin{tcolorbox}[annotbox]
The paper that founded matrix mechanics. Heisenberg replaces classical observables with
transition amplitudes and derives the commutation relations from a ``reinterpretation''
of kinematics. The historical origin of the operator formalism introduced in
Module~I.1 Lecture~L01.
\end{tcolorbox}
\end{bibentry}

\begin{bibentry}{t3light}
\textbf{Born, M.\ \& Jordan, P.}\\
\textit{Zur Quantenmechanik}\\
Z.~Phys.~\textbf{34}, 858--888 (1925). Translation in van der Waerden (ed.), 1967.\\[3pt]
\recommended\ [\textbf{P}]\\[3pt]
\begin{tcolorbox}[annotbox]
The first systematic formulation of matrix mechanics: canonical commutation relations
$[\hat{q},\hat{p}]=i\hbar$, energy matrix, equations of motion. The direct historical
predecessor of Module~I.3 Lecture~L07.
\end{tcolorbox}
\end{bibentry}

\begin{bibentry}{t3light}
\textbf{Born, M.}\\
\textit{Zur Quantenmechanik der Sto\ss vorg\"ange}\\
Z.~Phys.~\textbf{37}, 863--867 (1926). Translation in Wheeler \& Zurek,
\textit{Quantum Theory and Measurement}, Princeton, 1983.\\[3pt]
\essential\ [\textbf{P}]\\[3pt]
\begin{tcolorbox}[annotbox]
Two pages that changed physics: the Born rule $P=|\psi|^2$.
Must be read by all Tier~3 students. Directly relevant to Module~I.2 Lecture~L01
and the WRQ prompt ``How did Born's interpretation change the meaning of the wavefunction?''
\end{tcolorbox}
\end{bibentry}

\begin{bibentry}{t3light}
\textbf{Schr\"odinger, E.}\\
\textit{Quantisierung als Eigenwertproblem} (Parts I--IV)\\
Ann.~Phys.~\textbf{79}, 361 (1926); \textbf{79}, 489 (1926); \textbf{80}, 437 (1926);
\textbf{81}, 109 (1926). Translations in \textit{Wave Mechanics}, Pergamon, 1928.\\[3pt]
\recommended\ [\textbf{P}]\\[3pt]
\begin{tcolorbox}[annotbox]
The four papers establishing wave mechanics: the wave equation, its eigenvalue problem,
and the hydrogen atom solution. Part~II contains Schr\"odinger's proof that matrix
and wave mechanics are equivalent. Historical foundations for Module~I.3 Lectures~L05
and L11.
\end{tcolorbox}
\end{bibentry}

\begin{bibentry}{t3light}
\textbf{Heisenberg, W.}\\
\textit{\"Uber den anschaulichen Inhalt der quantentheoretischen Kinematik und
Mechanik}\\
Z.~Phys.~\textbf{43}, 172--198 (1927). Translation in Wheeler \& Zurek, 1983.\\[3pt]
\recommended\ [\textbf{P}]\\[3pt]
\begin{tcolorbox}[annotbox]
The uncertainty principle paper. Heisenberg derives the position-momentum uncertainty
relation via thought experiments (microscope). Historical foundation for Module~I.2
Lecture~L07 and Module~I.3 Lecture~L03.
\end{tcolorbox}
\end{bibentry}

\begin{bibentry}{t3light}
\textbf{Wigner, E.~P.}\\
\textit{On unitary representations of the inhomogeneous Lorentz group}\\
Ann.~Math.~\textbf{40}, 149--204 (1939).\\[3pt]
\useful\ [\textbf{P}]\\[3pt]
\begin{tcolorbox}[annotbox]
Original source for Wigner's theorem: symmetries preserve transition probabilities,
hence are implemented by unitary or antiunitary operators. Abstract accessible at
Tier~3; full proof requires Tier~5. Relevant to Module~I.2 Lectures~L01 and L03.
\end{tcolorbox}
\end{bibentry}

\begin{bibentry}{t3light}
\textbf{Einstein, A., Podolsky, B.\ \& Rosen, N.}\\
\textit{Can quantum-mechanical description of physical reality be considered complete?}\\
Phys.~Rev.~\textbf{47}, 777--780 (1935).\\[3pt]
\recommended\ [\textbf{P}]\\[3pt]
\begin{tcolorbox}[annotbox]
The EPR paper: the argument that quantum mechanics is incomplete via entangled states.
Three pages; fully accessible at Tier~3. Essential historical reading for Module~I.2
Lecture~L09 (tensor products and entanglement).
\end{tcolorbox}
\end{bibentry}

\begin{bibentry}{t3light}
\textbf{Bell, J.~S.}\\
\textit{On the Einstein Podolsky Rosen paradox}\\
Physics \textbf{1}, 195--200 (1964). Reprinted in Bell, \textit{Speakable and
Unspeakable in Quantum Mechanics}, Cambridge, 1987.\\[3pt]
\essential\ [\textbf{P}]\\[3pt]
\begin{tcolorbox}[annotbox]
Bell's theorem: no local hidden-variable theory can reproduce all predictions of quantum
mechanics. Six pages; the most important theorem in the foundations of quantum mechanics.
Essential Tier~3 reading; directly relevant to Module~I.2 Lecture~L07 and the ORQ on
operational axioms.
\end{tcolorbox}
\end{bibentry}

\newpage
%#############################################################
\section{Tier 4 --- Advanced Textbooks, Monographs \& Research Papers}
\label{sec:t4}
%#############################################################

\begin{tcolorbox}[tierheader=t4orange]
Tier 4 --- Master's (MSc): Advanced Textbooks, Specialised Topics \& Key Research Papers
\end{tcolorbox}

\medskip
\begin{tcolorbox}[keybox]
\textbf{Goal:} Command the full formalism; reproduce key derivations independently;
situate results in the broader research literature. Primary papers become the main
reading; textbooks serve as reference. Functional analysis and group theory background
assumed. Use alongside CPSG, SPJG, CPJG, and WRQ artifact types from the PSM.
\end{tcolorbox}

%------------------------------------------------------------
\subsection{Advanced Quantum Mechanics Textbooks}
%------------------------------------------------------------

\begin{bibentry}{t4light}
\textbf{Weinberg, S.}\\
\textit{Lectures on Quantum Mechanics} (2nd ed.)\\
Cambridge University Press, 2015.\\[3pt]
\essential\ [\textbf{T}]\\[3pt]
\begin{tcolorbox}[annotbox]
Covers foundations (postulates, measurement, entanglement), symmetries, and
perturbation theory at graduate level. Chapter~3 (``Mathematical Tools'') is the
cleanest graduate treatment of bra-ket including rigged Hilbert spaces.
Supports Module~I.2 Lectures~L11--L12 and Module~I.3 Lectures~L01--L09.
\end{tcolorbox}
\end{bibentry}

\begin{bibentry}{t4light}
\textbf{Galindo, A.\ \& Pascual, P.}\\
\textit{Quantum Mechanics} (2 vols.)\\
Springer, 1990--1991.\\[3pt]
\essential\ [\textbf{T}]\\[3pt]
\begin{tcolorbox}[annotbox]
The most mathematically complete standard graduate text. Volume~I develops Hilbert
space rigorously with spectral theory and self-adjoint operators; Volume~II covers
scattering, relativistic QM, and second quantisation. Recommended for Tier~4
students aiming at mathematical physics. Supports all of Series~I and~II.
\end{tcolorbox}
\end{bibentry}

\begin{bibentry}{t4light}
\textbf{Ballentine, L.~E.}\\
\textit{Quantum Mechanics: A Modern Development} (2nd ed.)\\
World Scientific, 2014.\\[3pt]
\essential\ [\textbf{T}]\\[3pt]
\begin{tcolorbox}[annotbox]
Statistical (ensemble) interpretation; detailed treatment of measurement, density
operators, and the relationship between classical and quantum mechanics. Excellent
for Module~I.2 Lectures~L07 and L11, and Module~I.3 Lecture~L10.
\end{tcolorbox}
\end{bibentry}

\begin{bibentry}{t4light}
\textbf{Cohen-Tannoudji, C., Diu, B.\ \& Lalo\"e, F.}\\
\textit{Quantum Mechanics} (2 vols., new ed.)\\
Wiley, 2019.\\[3pt]
\recommended\ [\textbf{T}]\\[3pt]
\begin{tcolorbox}[annotbox]
Encyclopaedic reference. The ``Complements'' to each chapter are especially valuable;
they cover applications, generalisations, and extensions that go beyond the main text.
Supports Modules~II.3 (angular momentum), II.6 (spectroscopy), and I.3.
\end{tcolorbox}
\end{bibentry}

\begin{bibentry}{t4light}
\textbf{Schwabl, F.}\\
\textit{Advanced Quantum Mechanics} (4th ed.)\\
Springer, 2008.\\[3pt]
\recommended\ [\textbf{T}]\\[3pt]
\begin{tcolorbox}[annotbox]
Second quantisation, many-body theory, relativistic quantum mechanics, and path
integrals at the MSc level. A compact bridge to Tier~5 monographs. Supports
Modules~II.2 and~II.4.
\end{tcolorbox}
\end{bibentry}

\begin{bibentry}{t4light}
\textbf{Le Bellac, M.}\\
\textit{Quantum Physics}\\
Cambridge University Press, 2006.\\[3pt]
\recommended\ [\textbf{T}]\\[3pt]
\begin{tcolorbox}[annotbox]
Modern graduate text emphasising quantum information, open systems, and quantum
optics alongside traditional topics. Excellent on density matrices, decoherence, and
quantum channels. Supports Module~I.2 Lectures~L08--L12 and Module~I.3 Lecture~L10.
\end{tcolorbox}
\end{bibentry}

%------------------------------------------------------------
\subsection{Symmetry, Group Theory \& Angular Momentum}
%------------------------------------------------------------

\begin{bibentry}{t4light}
\textbf{Wigner, E.~P.}\\
\textit{Group Theory and Its Application to the Quantum Mechanics of Atomic Spectra}\\
Academic Press, 1959 (translated from the 1931 German edition).\\[3pt]
\essential\ [\textbf{T/M}]\\[3pt]
\begin{tcolorbox}[annotbox]
The foundational text on symmetry in quantum mechanics. Wigner develops representation
theory, the rotation group, and selection rules. Chapters~1--15 support Module~II.3.
A primary source that remains the most authoritative treatment.
\end{tcolorbox}
\end{bibentry}

\begin{bibentry}{t4light}
\textbf{Georgi, H.}\\
\textit{Lie Algebras in Particle Physics} (2nd ed.)\\
Westview Press, 1999.\\[3pt]
\recommended\ [\textbf{T}]\\[3pt]
\begin{tcolorbox}[annotbox]
The standard physics graduate text on Lie algebras and representation theory.
$SU(2)$, $SU(3)$, and general semisimple algebras. Supports Module~II.3 Lectures~L01--L07
and the algebraic content of Module~I.1.
\end{tcolorbox}
\end{bibentry}

\begin{bibentry}{t4light}
\textbf{Biedenharn, L.~C.\ \& Louck, J.~D.}\\
\textit{Angular Momentum in Quantum Physics} (2 vols.)\\
Addison-Wesley, 1981. (Encyclopedia of Mathematics and Its Applications.)\\[3pt]
\recommended\ [\textbf{M}]\\[3pt]
\begin{tcolorbox}[annotbox]
The definitive encyclopaedic reference for angular momentum, Clebsch--Gordan
coefficients, Wigner symbols, and the Wigner--Eckart theorem. Supports
Module~II.3 Lectures~L05--L10 at graduate depth.
\end{tcolorbox}
\end{bibentry}

%------------------------------------------------------------
\subsection{Many-Body Theory}
%------------------------------------------------------------

\begin{bibentry}{t4light}
\textbf{Fetter, A.~L.\ \& Walecka, J.~D.}\\
\textit{Quantum Theory of Many-Particle Systems}\\
McGraw-Hill, 1971 (Dover reprint 2003).\\[3pt]
\essential\ [\textbf{T}]\\[3pt]
\begin{tcolorbox}[annotbox]
The standard graduate text for Series~II Modules~II.4--II.5. Second quantisation,
Green's functions, Feynman diagrams, and linear response developed systematically.
Chapter~1 pairs with Module~II.4 Lecture~L01; Chapters~7--9 with Module~II.5.
\end{tcolorbox}
\end{bibentry}

\begin{bibentry}{t4light}
\textbf{Bruus, H.\ \& Flensberg, K.}\\
\textit{Introduction to Many-Body Quantum Theory in Condensed Matter Physics}\\
Oxford University Press, 2004. Freely available online.\\[3pt]
\essential\ [\textbf{T}]\\[3pt]
\begin{tcolorbox}[annotbox]
More accessible than Fetter--Walecka. Excellent on second quantisation, diagrammatics,
and Green's functions for condensed matter applications. Primary reading for
Module~II.4 Lectures~L01--L07.
\end{tcolorbox}
\end{bibentry}

\begin{bibentry}{t4light}
\textbf{Stefanucci, G.\ \& van Leeuwen, R.}\\
\textit{Nonequilibrium Many-Body Theory of Quantum Systems}\\
Cambridge University Press, 2013.\\[3pt]
\essential\ [\textbf{T}]\\[3pt]
\begin{tcolorbox}[annotbox]
The definitive text for Module~II.5 (NEGF). Derives the Keldysh formalism
rigorously from first principles; self-energies, Kadanoff--Baym equations,
and steady-state transport developed in full. Chapters~4--8 are primary reading
for Module~II.5 Lectures~L02--L09.
\end{tcolorbox}
\end{bibentry}

\begin{bibentry}{t4light}
\textbf{Mahan, G.~D.}\\
\textit{Many-Particle Physics} (3rd ed.)\\
Kluwer/Plenum, 2000.\\[3pt]
\recommended\ [\textbf{T}]\\[3pt]
\begin{tcolorbox}[annotbox]
A comprehensive graduate text on many-body methods in condensed matter. Green's
functions, self-energy, diagrammatic perturbation theory, and transport.
Supports Modules~II.4--II.6 as a reference text alongside Fetter--Walecka.
\end{tcolorbox}
\end{bibentry}

%------------------------------------------------------------
\subsection{Quantum Statistical Mechanics}
%------------------------------------------------------------

\begin{bibentry}{t4light}
\textbf{Huang, K.}\\
\textit{Statistical Mechanics} (2nd ed.)\\
Wiley, 1987.\\[3pt]
\recommended\ [\textbf{T}]\\[3pt]
\begin{tcolorbox}[annotbox]
The standard graduate statistical mechanics text. Chapters on quantum statistical
mechanics, Bose--Einstein condensation, and the Ising model support Module~II.7.
\end{tcolorbox}
\end{bibentry}

\begin{bibentry}{t4light}
\textbf{Sachdev, S.}\\
\textit{Quantum Phase Transitions} (2nd ed.)\\
Cambridge University Press, 2011.\\[3pt]
\recommended\ [\textbf{T}]\\[3pt]
\begin{tcolorbox}[annotbox]
The authoritative text on quantum phase transitions, quantum criticality, and
non-Fermi liquids. Supports Module~II.7 Lectures~L09--L12.
\end{tcolorbox}
\end{bibentry}

%------------------------------------------------------------
\subsection{Quantum Chemistry (Series III)}
%------------------------------------------------------------

\begin{bibentry}{t4light}
\textbf{Szabo, A.\ \& Ostlund, N.~S.}\\
\textit{Modern Quantum Chemistry: Introduction to Advanced Electronic Structure Theory}\\
McGraw-Hill, 1989 (Dover reprint 1996).\\[3pt]
\essential\ [\textbf{T}]\\[3pt]
\begin{tcolorbox}[annotbox]
The standard entry-level graduate text for Series~III. Hartree--Fock, CI, MP2,
and the rudiments of DFT developed concisely with full derivations. Chapters~1--3
support Module~III.1; Chapters~4--6 support Module~III.2 Lectures~L05--L11.
\end{tcolorbox}
\end{bibentry}

\begin{bibentry}{t4light}
\textbf{Helgaker, T., J\o rgensen, P.\ \& Olsen, J.}\\
\textit{Molecular Electronic-Structure Theory}\\
Wiley, 2000.\\[3pt]
\essential\ [\textbf{M}]\\[3pt]
\begin{tcolorbox}[annotbox]
The definitive graduate reference for Modules~III.2--III.3. 900 pages covering
basis sets, HF, MP, CC, MCSCF, and response theory. Use as the reference text
alongside Module~III.2 from Lecture~L07 onward.
\end{tcolorbox}
\end{bibentry}

\begin{bibentry}{t4light}
\textbf{Koch, W.\ \& Holthausen, M.~C.}\\
\textit{A Chemist's Guide to Density Functional Theory} (2nd ed.)\\
Wiley-VCH, 2001.\\[3pt]
\recommended\ [\textbf{T}]\\[3pt]
\begin{tcolorbox}[annotbox]
Accessible DFT text for Tier~4 chemistry students. Good companion for
Module~III.3 Lectures~L01--L05.
\end{tcolorbox}
\end{bibentry}

\begin{bibentry}{t4light}
\textbf{Martin, R.~M.}\\
\textit{Electronic Structure: Basic Theory and Practical Methods}\\
Cambridge University Press, 2004.\\[3pt]
\recommended\ [\textbf{T}]\\[3pt]
\begin{tcolorbox}[annotbox]
DFT from first principles: plane waves, pseudopotentials, PAW, and band structure.
Complements the molecular orbital perspective of Module~III.2 with the solid-state
perspective of Module~III.2 Lecture~L16 and Module~III.3.
\end{tcolorbox}
\end{bibentry}

%------------------------------------------------------------
\subsection{Path Integrals}
%------------------------------------------------------------

\begin{bibentry}{t4light}
\textbf{Feynman, R.~P.\ \& Hibbs, A.~R.}\\
\textit{Quantum Mechanics and Path Integrals} (emended ed., Styer)\\
Dover, 2010 (original McGraw-Hill, 1965).\\[3pt]
\essential\ [\textbf{T}]\\[3pt]
\begin{tcolorbox}[annotbox]
The original path-integral textbook. Develops the Feynman sum-over-histories
from physical principles. Essential reading for Module~II.2 Lectures~L01--L04.
\end{tcolorbox}
\end{bibentry}

\begin{bibentry}{t4light}
\textbf{Zinn-Justin, J.}\\
\textit{Quantum Field Theory and Critical Phenomena} (4th ed.)\\
Oxford University Press, 2002.\\[3pt]
\recommended\ [\textbf{M}]\\[3pt]
\begin{tcolorbox}[annotbox]
The standard graduate reference for path integrals in field theory. Parts~I--II
(functional integrals, Gaussian integrals, perturbation theory) support
Module~II.2 Lectures~L03--L09.
\end{tcolorbox}
\end{bibentry}

%------------------------------------------------------------
\subsection{Key Tier 4 Research Papers (QM \& QC)}
%------------------------------------------------------------

\begin{bibentry}{t4light}
\textbf{Hohenberg, P.\ \& Kohn, W.}\\
\textit{Inhomogeneous electron gas}\\
Phys.~Rev.~\textbf{136}, B864--B871 (1964).\\[3pt]
\essential\ [\textbf{P}]\\[3pt]
\begin{tcolorbox}[annotbox]
Proves the two Hohenberg--Kohn theorems establishing DFT. Primary reading for
Module~III.3 Lecture~L01.
\end{tcolorbox}
\end{bibentry}

\begin{bibentry}{t4light}
\textbf{Kohn, W.\ \& Sham, L.~J.}\\
\textit{Self-consistent equations including exchange and correlation effects}\\
Phys.~Rev.~\textbf{140}, A1133--A1138 (1965).\\[3pt]
\essential\ [\textbf{P}]\\[3pt]
\begin{tcolorbox}[annotbox]
The Kohn--Sham DFT paper: maps the interacting problem onto a non-interacting
Kohn--Sham system. Primary reading for Module~III.3 Lecture~L03.
\end{tcolorbox}
\end{bibentry}

\begin{bibentry}{t4light}
\textbf{Perdew, J.~P.\ \& Zunger, A.}\\
\textit{Self-interaction correction to density-functional approximations for
many-electron systems}\\
Phys.~Rev.~B \textbf{23}, 5048 (1981).\\[3pt]
\recommended\ [\textbf{P}]\\[3pt]
\begin{tcolorbox}[annotbox]
Foundational DFT paper on self-interaction error and LDA. Key reading for
Module~III.3 Lecture~L04 (XC functionals).
\end{tcolorbox}
\end{bibentry}

\begin{bibentry}{t4light}
\textbf{Becke, A.~D.}\\
\textit{Density-functional thermochemistry. III. The role of exact exchange}\\
J.~Chem.~Phys.~\textbf{98}, 5648--5652 (1993).\\[3pt]
\recommended\ [\textbf{P}]\\[3pt]
\begin{tcolorbox}[annotbox]
The paper introducing the B3LYP hybrid functional --- the most widely used
functional in quantum chemistry. Essential context for Module~III.3 Lecture~L04.
\end{tcolorbox}
\end{bibentry}

\begin{bibentry}{t4light}
\textbf{Perdew, J.~P., Burke, K.\ \& Ernzerhof, M.}\\
\textit{Generalized gradient approximation made simple}\\
Phys.~Rev.~Lett.~\textbf{77}, 3865--3868 (1996).\\[3pt]
\recommended\ [\textbf{P}]\\[3pt]
\begin{tcolorbox}[annotbox]
The PBE functional paper --- the standard GGA functional in solid-state calculations.
Essential alongside B3LYP for Module~III.3 Lecture~L04.
\end{tcolorbox}
\end{bibentry}

\begin{bibentry}{t4light}
\textbf{Kadanoff, L.~P.\ \& Baym, G.}\\
\textit{Quantum Statistical Mechanics}\\
Benjamin, 1962 (Perseus reprint 1989).\\[3pt]
\essential\ [\textbf{M}]\\[3pt]
\begin{tcolorbox}[annotbox]
Original derivation of the Kadanoff--Baym equations for the two-time Green's function.
Primary reading for Module~II.5 Lecture~L05.
\end{tcolorbox}
\end{bibentry}

\begin{bibentry}{t4light}
\textbf{Berry, M.~V.}\\
\textit{Quantal phase factors accompanying adiabatic changes}\\
Proc.~R.~Soc.~Lond.~A \textbf{392}, 45--57 (1984).\\[3pt]
\essential\ [\textbf{P}]\\[3pt]
\begin{tcolorbox}[annotbox]
Berry's original paper on geometric phases. Directly relevant to Module~II.2
Lecture~L12 (Berry phase from path integrals) and Module~I.3 Lecture~L06
(non-commuting time evolution).
\end{tcolorbox}
\end{bibentry}

\begin{bibentry}{t4light}
\textbf{Pulay, P.}\\
\textit{Ab initio calculation of force constants and equilibrium geometries
in polyatomic molecules}\\
Mol.~Phys.~\textbf{17}, 197--204 (1969).\\[3pt]
\recommended\ [\textbf{P}]\\[3pt]
\begin{tcolorbox}[annotbox]
Introduces analytic gradient theory (``Pulay forces'') in quantum chemistry ---
the foundation of modern geometry optimisation. Supports Module~III.1 Lecture~L03
(Feynman--Hellmann theorem and force calculations).
\end{tcolorbox}
\end{bibentry}

\begin{bibentry}{t4light}
\textbf{Runge, E.\ \& Gross, E.~K.~U.}\\
\textit{Density-functional theory for time-dependent systems}\\
Phys.~Rev.~Lett.~\textbf{52}, 997--1000 (1984).\\[3pt]
\recommended\ [\textbf{P}]\\[3pt]
\begin{tcolorbox}[annotbox]
The Runge--Gross theorem establishing time-dependent DFT. Background reading for
Module~III.3 and for the linear response content of Module~II.6.
\end{tcolorbox}
\end{bibentry}

\begin{bibentry}{t4light}
\textbf{Kubo, R.}\\
\textit{Statistical-mechanical theory of irreversible processes. I.}\\
J.~Phys.~Soc.~Japan \textbf{12}, 570--586 (1957).\\[3pt]
\essential\ [\textbf{P}]\\[3pt]
\begin{tcolorbox}[annotbox]
The Kubo formula for linear response. Primary reading for Module~II.6 Lecture~L03.
\end{tcolorbox}
\end{bibentry}

\newpage
%#############################################################
\section{Tier 5 --- Research Monographs, Specialised Topics \& Primary Literature}
\label{sec:t5}
%#############################################################

\begin{tcolorbox}[tierheader=t5red]
Tier 5 --- PhD: Research Monographs, Specialised Topics \& Full Primary Literature
\end{tcolorbox}

\medskip
\begin{tcolorbox}[keybox]
\textbf{Goal:} Achieve research-level mastery; read primary literature directly;
identify open problems; connect the programme to current research frontiers.
Primary papers and monographs are the main reading. Use alongside CPJG and ORQ
artifact types from the PSM.
\end{tcolorbox}

%------------------------------------------------------------
\subsection{Functional Analysis \& Mathematical Foundations of QM}
%------------------------------------------------------------

\begin{bibentry}{t5light}
\textbf{Reed, M.\ \& Simon, B.}\\
\textit{Methods of Modern Mathematical Physics} (4 vols.)\\
Academic Press, 1972--1979.\\
Vol.~I: \textit{Functional Analysis}. Vol.~II: \textit{Fourier Analysis, Self-Adjointness}.
Vol.~III: \textit{Scattering Theory}. Vol.~IV: \textit{Analysis of Operators}.\\[3pt]
\essential\ [\textbf{M}]\\[3pt]
\begin{tcolorbox}[annotbox]
The mathematical foundation of quantum mechanics at the PhD level.
Vol.~I: Hilbert spaces, operators, spectral theory in full rigour.
Vol.~II: self-adjoint extensions, deficiency indices, the Gel'fand triple.
Vol.~IV: spectral theory of Schr\"odinger operators.
Essential for Module~I.2 Lecture~L12 and Module~I.3 Lecture~L04.
Every PhD student in mathematical physics should work through Vol.~I--II.
\end{tcolorbox}
\end{bibentry}

\begin{bibentry}{t5light}
\textbf{Rudin, W.}\\
\textit{Functional Analysis} (2nd ed.)\\
McGraw-Hill, 1991.\\[3pt]
\essential\ [\textbf{M}]\\[3pt]
\begin{tcolorbox}[annotbox]
Banach and Hilbert spaces, the spectral theorem for bounded operators, distributions.
Chapter~12 (unbounded operators) is directly relevant to the self-adjoint extension
problem for $\hat{p}$. Supports Module~I.2 Lecture~L12 at PhD level.
\end{tcolorbox}
\end{bibentry}

\begin{bibentry}{t5light}
\textbf{Blank, J., Exner, P.\ \& Havl\'{\i}\v{c}ek, M.}\\
\textit{Hilbert Space Operators in Quantum Physics} (2nd ed.)\\
Springer, 2008.\\[3pt]
\essential\ [\textbf{M}]\\[3pt]
\begin{tcolorbox}[annotbox]
The most physics-oriented rigorous mathematical treatment. Spectral measures,
rigged Hilbert spaces, quantum dynamics, and quantum fields. Directly supports
Module~I.2 Lectures~L11--L12 and Module~I.3 Lectures~L04--L08.
\end{tcolorbox}
\end{bibentry}

\begin{bibentry}{t5light}
\textbf{Gel'fand, I.~M.\ \& Vilenkin, N.~Ya.}\\
\textit{Generalized Functions, Vol.~4: Applications of Harmonic Analysis}\\
Academic Press, 1964.\\[3pt]
\recommended\ [\textbf{M}]\\[3pt]
\begin{tcolorbox}[annotbox]
Original source for the rigged Hilbert space (Gel'fand triple) framework
$\Phi\subset\mathcal{H}\subset\Phi'$. Chapter~1 is the foundational reference for
Module~I.2 Lecture~L12.
\end{tcolorbox}
\end{bibentry}

\begin{bibentry}{t5light}
\textbf{Thirring, W.}\\
\textit{A Course in Mathematical Physics} (4 vols.)\\
Springer, 1978--1983.\\
Vol.~3: \textit{Quantum Mechanics of Atoms and Molecules}. Vol.~4: \textit{Quantum Mechanics of Large Systems}.\\[3pt]
\recommended\ [\textbf{M}]\\[3pt]
\begin{tcolorbox}[annotbox]
Mathematically rigorous treatment of quantum mechanics using $C^*$-algebras and
operator algebras. Vol.~3 covers the spectral theory of atomic Hamiltonians;
Vol.~4 treats quantum statistical mechanics and stability of matter. Essential
background for algebraic approaches to Module~II.7.
\end{tcolorbox}
\end{bibentry}

\begin{bibentry}{t5light}
\textbf{Lieb, E.~H.\ \& Loss, M.}\\
\textit{Analysis} (2nd ed.)\\
American Mathematical Society, 2001.\\[3pt]
\recommended\ [\textbf{M}]\\[3pt]
\begin{tcolorbox}[annotbox]
Covers the functional analysis and real analysis needed for rigorous quantum mechanics:
$L^p$ spaces, Sobolev spaces, variational methods, and sharp inequalities. The
Lieb--Thirring inequality and stability-of-matter results are treated here. PhD-level
background for Modules~III.1 and~III.3.
\end{tcolorbox}
\end{bibentry}

%------------------------------------------------------------
\subsection{Foundations of Quantum Mechanics}
%------------------------------------------------------------

\begin{bibentry}{t5light}
\textbf{von Neumann, J.}\\
\textit{Mathematical Foundations of Quantum Mechanics}\\
Princeton University Press, 1955 (translated from German, 1932).\\[3pt]
\essential\ [\textbf{M}]\\[3pt]
\begin{tcolorbox}[annotbox]
The foundational monograph. Chapters~I--II: Hilbert space and operators.
Chapter~III: quantum statistics and measurement. Chapter~VI: the measurement
problem. Chapters~III--VI require Tier~5 preparation. The historical and
conceptual foundation for Module~I.2 Lecture~L12 and Module~I.3 Lecture~L10.
\end{tcolorbox}
\end{bibentry}

\begin{bibentry}{t5light}
\textbf{Haag, R.}\\
\textit{Local Quantum Physics: Fields, Particles, Algebras} (2nd ed.)\\
Springer, 1996.\\[3pt]
\recommended\ [\textbf{M}]\\[3pt]
\begin{tcolorbox}[annotbox]
Algebraic QFT (AQFT): observables as $C^*$-algebras, superselection sectors,
Tomita--Takesaki modular theory. The deepest mathematical treatment of locality
and the measurement problem. Background for Series~II Module~II.4 at research level.
\end{tcolorbox}
\end{bibentry}

\begin{bibentry}{t5light}
\textbf{Kraus, K.}\\
\textit{States, Effects, and Operations: Fundamental Notions of Quantum Theory}\\
Springer Lecture Notes in Physics \textbf{190}, 1983.\\[3pt]
\recommended\ [\textbf{M}]\\[3pt]
\begin{tcolorbox}[annotbox]
Kraus operators, CPTP maps, and the operational (instrument) formulation of
quantum mechanics. The mathematical background for Module~I.2 Lecture~L11 and
Module~I.3 Lecture~L10 at PhD depth.
\end{tcolorbox}
\end{bibentry}

\begin{bibentry}{t5light}
\textbf{Busch, P., Lahti, P.\ \& Mittelstaedt, P.}\\
\textit{The Quantum Theory of Measurement} (2nd ed.)\\
Springer Lecture Notes in Physics \textbf{m2}, 1996.\\[3pt]
\recommended\ [\textbf{M}]\\[3pt]
\begin{tcolorbox}[annotbox]
Rigorous axiomatic treatment of quantum measurement: POVMs, instruments, repeatability,
and the L\"uders rule. The mathematical foundation for Module~I.2 Lecture~L07 and L11
at PhD level.
\end{tcolorbox}
\end{bibentry}

%------------------------------------------------------------
\subsection{Advanced Many-Body Theory \& Quantum Field Theory}
%------------------------------------------------------------

\begin{bibentry}{t5light}
\textbf{Negele, J.~W.\ \& Orland, H.}\\
\textit{Quantum Many-Particle Systems}\\
Addison-Wesley, 1988 (Perseus reprint 1998).\\[3pt]
\essential\ [\textbf{M}]\\[3pt]
\begin{tcolorbox}[annotbox]
Path integrals and Green's functions at research level: functional methods,
coherent state path integrals, finite-temperature diagrammatics. Supports
Modules~II.2, II.4, and II.5.
\end{tcolorbox}
\end{bibentry}

\begin{bibentry}{t5light}
\textbf{Peskin, M.~E.\ \& Schroeder, D.~V.}\\
\textit{An Introduction to Quantum Field Theory}\\
Addison-Wesley, 1995.\\[3pt]
\recommended\ [\textbf{T}]\\[3pt]
\begin{tcolorbox}[annotbox]
The standard QFT graduate text. Parts~I--II (canonical quantisation, Feynman rules,
one-loop divergences) provide the QFT background for Module~II.4 Lecture~L12
(beyond mean field, connection to QFT) and Module~II.2 Lecture~L09.
\end{tcolorbox}
\end{bibentry}

\begin{bibentry}{t5light}
\textbf{Altland, A.\ \& Simons, B.~D.}\\
\textit{Condensed Matter Field Theory} (2nd ed.)\\
Cambridge University Press, 2010.\\[3pt]
\essential\ [\textbf{T}]\\[3pt]
\begin{tcolorbox}[annotbox]
The modern graduate text bridging quantum field theory and condensed matter. Path
integrals, functional methods, the renormalisation group, and topology in condensed
matter. Primary reading for Module~II.2 and a key reference for Modules~II.4--II.5.
\end{tcolorbox}
\end{bibentry}

\begin{bibentry}{t5light}
\textbf{Abrikosov, A.~A., Gorkov, L.~P.\ \& Dzyaloshinski, I.~E.}\\
\textit{Methods of Quantum Field Theory in Statistical Physics}\\
Dover, 1975 (translated from the 1962 Russian edition).\\[3pt]
\recommended\ [\textbf{M}]\\[3pt]
\begin{tcolorbox}[annotbox]
The classic Soviet text on diagrammatic techniques for many-body systems. Green's
functions, self-energy, phonon--electron interactions, and superconductivity.
Supports Modules~II.4 and II.5 at research level.
\end{tcolorbox}
\end{bibentry}

%------------------------------------------------------------
\subsection{Quantum Information \& Open Quantum Systems}
%------------------------------------------------------------

\begin{bibentry}{t5light}
\textbf{Nielsen, M.~A.\ \& Chuang, I.~L.}\\
\textit{Quantum Computation and Quantum Information}\\
Cambridge University Press, 2000.\\[3pt]
\essential\ [\textbf{M}]\\[3pt]
\begin{tcolorbox}[annotbox]
The definitive reference for quantum information and quantum computing. Chapters~2
(linear algebra and density operators), 8 (quantum noise), and 9--10 (error
correction) are directly relevant to Module~I.2 Lectures~L09--L11 and Module~I.3
Lecture~L10 at research depth.
\end{tcolorbox}
\end{bibentry}

\begin{bibentry}{t5light}
\textbf{Breuer, H.-P.\ \& Petruccione, F.}\\
\textit{The Theory of Open Quantum Systems}\\
Oxford University Press, 2002.\\[3pt]
\essential\ [\textbf{M}]\\[3pt]
\begin{tcolorbox}[annotbox]
The standard monograph on open quantum systems: Lindblad master equations,
quantum Brownian motion, decoherence, and quantum trajectories. Primary reference
for Module~I.3 Lecture~L10 and Module~II.5 Lecture~L11.
\end{tcolorbox}
\end{bibentry}

\begin{bibentry}{t5light}
\textbf{Carmichael, H.~J.}\\
\textit{An Open Systems Approach to Quantum Optics}\\
Springer Lecture Notes in Physics \textbf{m18}, 1993.\\[3pt]
\recommended\ [\textbf{M}]\\[3pt]
\begin{tcolorbox}[annotbox]
Quantum trajectories and the stochastic Schr\"odinger equation. A key tool for
numerical simulation of open quantum systems. Supports Module~I.3 Lecture~L10
and Module~II.6 at research depth.
\end{tcolorbox}
\end{bibentry}

%------------------------------------------------------------
\subsection{Advanced Quantum Chemistry (Series III Research Level)}
%------------------------------------------------------------

\begin{bibentry}{t5light}
\textbf{Martin, R.~M., Reining, L.\ \& Ceperley, D.}\\
\textit{Interacting Electrons: Theory and Computational Approaches}\\
Cambridge University Press, 2016.\\[3pt]
\essential\ [\textbf{M}]\\[3pt]
\begin{tcolorbox}[annotbox]
The definitive modern text connecting DFT, many-body perturbation theory ($GW$, BSE),
and quantum Monte Carlo. All of Series~III covered at research depth.
\end{tcolorbox}
\end{bibentry}

\begin{bibentry}{t5light}
\textbf{Parr, R.~G.\ \& Yang, W.}\\
\textit{Density-Functional Theory of Atoms and Molecules}\\
Oxford University Press, 1989.\\[3pt]
\essential\ [\textbf{M}]\\[3pt]
\begin{tcolorbox}[annotbox]
The foundational DFT monograph by its principal architects. Develops HK theorems,
KS equations, the local density approximation, and chemical reactivity indices
(Fukui functions, hardness/softness). Primary reference for Module~III.3 Lectures~L01--L05.
\end{tcolorbox}
\end{bibentry}

\begin{bibentry}{t5light}
\textbf{Dreizler, R.~M.\ \& Gross, E.~K.~U.}\\
\textit{Density Functional Theory: An Approach to the Quantum Many-Body Problem}\\
Springer, 1990.\\[3pt]
\recommended\ [\textbf{M}]\\[3pt]
\begin{tcolorbox}[annotbox]
Rigorous treatment of DFT from a many-body physics perspective. Covers HK theorems,
KS equations, LDA, and TDDFT. Supports Module~III.3 at research level.
\end{tcolorbox}
\end{bibentry}

%------------------------------------------------------------
\subsection{Essential PhD-Level Primary Papers}
%------------------------------------------------------------

\begin{bibentry}{t5light}
\textbf{Stone, M.~H.}\\
\textit{On one-parameter unitary groups in Hilbert space}\\
Ann.~Math.~\textbf{33}, 643--648 (1932).\\[3pt]
\essential\ [\textbf{P}]\\[3pt]
\begin{tcolorbox}[annotbox]
Stone's theorem: the mathematical basis for $\hat{U}(t)=e^{-i\hat{H}t/\hbar}$.
Primary reading for Module~I.3 Lecture~L04 at PhD level.
\end{tcolorbox}
\end{bibentry}

\begin{bibentry}{t5light}
\textbf{Gleason, A.~M.}\\
\textit{Measures on the closed subspaces of a Hilbert space}\\
J.~Math.~Mech.~\textbf{6}, 885--893 (1957).\\[3pt]
\essential\ [\textbf{P}]\\[3pt]
\begin{tcolorbox}[annotbox]
Derives the Born rule from measure theory on the lattice of closed subspaces:
the deepest axiomatic justification of $P=|\langle a|\psi\rangle|^2$.
Essential for Module~I.2 Lecture~L07 at PhD level.
\end{tcolorbox}
\end{bibentry}

\begin{bibentry}{t5light}
\textbf{Lindblad, G.}\\
\textit{On the generators of quantum dynamical semigroups}\\
Commun.~Math.~Phys.~\textbf{48}, 119--130 (1976).\\[3pt]
\essential\ [\textbf{P}]\\[3pt]
\begin{tcolorbox}[annotbox]
Derives the most general Markovian master equation (Lindblad equation) from
complete positivity and trace preservation. Primary reading for Module~I.3
Lecture~L10 and Module~II.5 Lecture~L11.
\end{tcolorbox}
\end{bibentry}

\begin{bibentry}{t5light}
\textbf{Gorini, V., Kossakowski, A.\ \& Sudarshan, E.~C.~G.}\\
\textit{Completely positive dynamical semigroups of N-level systems}\\
J.~Math.~Phys.~\textbf{17}, 821--825 (1976).\\[3pt]
\essential\ [\textbf{P}]\\[3pt]
\begin{tcolorbox}[annotbox]
Simultaneous with Lindblad (1976): independently derives the GKSL master equation
for finite-dimensional systems. Together with Lindblad's paper, establishes the
complete mathematical framework for open quantum dynamics.
\end{tcolorbox}
\end{bibentry}

\begin{bibentry}{t5light}
\textbf{Aspect, A., Grangier, P.\ \& Roger, G.}\\
\textit{Experimental tests of Bell's inequalities using time-varying analyzers}\\
Phys.~Rev.~Lett.~\textbf{49}, 1804--1807 (1982).\\[3pt]
\essential\ [\textbf{P}]\\[3pt]
\begin{tcolorbox}[annotbox]
The definitive experimental test of Bell's theorem. Closes the locality loophole
and firmly establishes quantum nonlocality. Essential reading alongside Bell (1964)
for Module~I.2 Lecture~L07 and L09 at PhD level.
\end{tcolorbox}
\end{bibentry}

\begin{bibentry}{t5light}
\textbf{Heisenberg, W.}\\
\textit{Zur Theorie des Ferromagnetismus}\\
Z.~Phys.~\textbf{49}, 619--636 (1928).\\[3pt]
\recommended\ [\textbf{P}]\\[3pt]
\begin{tcolorbox}[annotbox]
The Heisenberg exchange interaction: exchange integral from antisymmetry of
electron wavefunctions. Historical foundation for Module~II.4 Lecture~L11
(quantum magnetism).
\end{tcolorbox}
\end{bibentry}

\begin{bibentry}{t5light}
\textbf{Bardeen, J., Cooper, L.~N.\ \& Schrieffer, J.~R.}\\
\textit{Theory of superconductivity}\\
Phys.~Rev.~\textbf{108}, 1175--1204 (1957).\\[3pt]
\essential\ [\textbf{P}]\\[3pt]
\begin{tcolorbox}[annotbox]
The BCS theory: Cooper pairs, the BCS gap equation, and the Bogoliubov--Valatin
transformation. Primary reading for Module~II.4 Lecture~L09 (BCS superconductivity).
\end{tcolorbox}
\end{bibentry}

\begin{bibentry}{t5light}
\textbf{Keldysh, L.~V.}\\
\textit{Diagram technique for nonequilibrium processes}\\
Zh.~Eksp.~Teor.~Fiz.~\textbf{47}, 1515 (1964); Sov.~Phys.~JETP \textbf{20}, 1018 (1965).\\[3pt]
\essential\ [\textbf{P}]\\[3pt]
\begin{tcolorbox}[annotbox]
The Keldysh contour and closed-time-path formalism for non-equilibrium Green's
functions. The foundational paper for Module~II.5 Lectures~L02--L04.
\end{tcolorbox}
\end{bibentry}

\begin{bibentry}{t5light}
\textbf{Meir, Y.\ \& Wingreen, N.~S.}\\
\textit{Landauer formula for the current through an interacting electron region}\\
Phys.~Rev.~Lett.~\textbf{68}, 2512--2515 (1992).\\[3pt]
\recommended\ [\textbf{P}]\\[3pt]
\begin{tcolorbox}[annotbox]
Derives the Meir--Wingreen formula for current through a correlated region.
The central result of Module~II.5 Lecture~L08 (steady-state transport).
\end{tcolorbox}
\end{bibentry}

\begin{bibentry}{t5light}
\textbf{Aharonov, Y.\ \& Bohm, D.}\\
\textit{Significance of electromagnetic potentials in the quantum theory}\\
Phys.~Rev.~\textbf{115}, 485--491 (1959).\\[3pt]
\essential\ [\textbf{P}]\\[3pt]
\begin{tcolorbox}[annotbox]
The Aharonov--Bohm effect: the vector potential is physically significant even
in regions of zero field. The quantum mechanics of gauge fields. Primary reading
for Module~II.2 Lecture~L10.
\end{tcolorbox}
\end{bibentry}

\begin{bibentry}{t5light}
\textbf{Zurek, W.~H.}\\
\textit{Pointer basis of quantum apparatus: Into what mixture does the wave packet collapse?}\\
Phys.~Rev.~D \textbf{24}, 1516--1525 (1981).\\[3pt]
\recommended\ [\textbf{P}]\\[3pt]
\begin{tcolorbox}[annotbox]
Introduces the concept of environment-induced superselection (``einselection'') and
pointer states. The foundational paper of decoherence theory. Primary reading for
Module~I.3 Lecture~L10 (open systems and measurement).
\end{tcolorbox}
\end{bibentry}

\begin{bibentry}{t5light}
\textbf{Dirac, P.~A.~M.}\\
\textit{The quantum theory of the emission and absorption of radiation}\\
Proc.~R.~Soc.~Lond.~A \textbf{114}, 243--265 (1927).\\[3pt]
\recommended\ [\textbf{P}]\\[3pt]
\begin{tcolorbox}[annotbox]
The first paper on quantum electrodynamics: quantisation of the electromagnetic
field, creation/annihilation operators, and spontaneous emission. Historical
foundation for Module~II.4 and Module~II.6 Lecture~L01.
\end{tcolorbox}
\end{bibentry}

\begin{bibentry}{t5light}
\textbf{Wigner, E.~P.}\\
\textit{The unreasonable effectiveness of mathematics in the natural sciences}\\
Commun.~Pure Appl.~Math.~\textbf{13}, 1--14 (1960).\\[3pt]
\useful\ [\textbf{P}]\\[3pt]
\begin{tcolorbox}[annotbox]
A philosophical essay on the mysterious alignment between abstract mathematics and
physical law. Stimulating reading for PhD students reflecting on the foundations of
the Hilbert space formalism.
\end{tcolorbox}
\end{bibentry}

%------------------------------------------------------------
\subsection{Review Articles \& Survey Monographs}
%------------------------------------------------------------

\begin{bibentry}{t5light}
\textbf{Leggett, A.~J.}\\
\textit{Quantum Liquids: Bose Condensation and Cooper Pairing in Condensed-Matter Systems}\\
Oxford University Press, 2006. [\textbf{M}]\\[3pt]
\essential\\[3pt]
\begin{tcolorbox}[annotbox]
Research-level survey of BEC and BCS from a unified perspective. Supports
Module~II.4 Lectures~L08--L09 and connects to Module~II.7.
\end{tcolorbox}
\end{bibentry}

\begin{bibentry}{t5light}
\textbf{Zurek, W.~H.}\ [\textbf{R}]\\
\textit{Decoherence, einselection, and the quantum origins of the classical}\\
Rev.~Mod.~Phys.~\textbf{75}, 715--775 (2003).\\[3pt]
\essential\\[3pt]
\begin{tcolorbox}[annotbox]
The comprehensive review of decoherence theory. Covers environment-induced
superselection, pointer states, quantum Darwinism, and the emergence of classicality.
Primary review reading for Module~I.3 Lecture~L10 and Module~II.5 Lecture~L11.
\end{tcolorbox}
\end{bibentry}

\begin{bibentry}{t5light}
\textbf{Koch, J.\ et al.}\ [\textbf{R}]\\
\textit{Charge-insensitive qubit design derived from the Cooper pair box}\\
Phys.~Rev.~A \textbf{76}, 042319 (2007).\\[3pt]
\recommended\\[3pt]
\begin{tcolorbox}[annotbox]
The transmon qubit paper. A key paper in the application of Module~I.2's formalism
(two-level systems, Josephson junctions) to superconducting quantum computing.
Excellent PhD-level example of the Jaynes--Cummings Hamiltonian in practice.
\end{tcolorbox}
\end{bibentry}

\begin{bibentry}{t5light}
\textbf{Rotter, S.\ \& Fan, S.}\ [\textbf{R}]\\
\textit{Optical cavities during the age of quantum information: Cavity QED and non-Hermitian physics}\\
Science \textbf{358}, 1444--1447 (2017).\\[3pt]
\useful\\[3pt]
\begin{tcolorbox}[annotbox]
Modern review connecting cavity QED (Module~II.6 Lecture~L12) to non-Hermitian
physics and exceptional points. Represents a current research frontier.
\end{tcolorbox}
\end{bibentry}

\newpage
%#############################################################
\section{Full Alphabetical Catalogue}
\label{sec:fullcat}
%#############################################################

\begin{tcolorbox}[keybox]
Complete alphabetically ordered catalogue grouped by category.
Each entry carries tier badge(s), type tag, and star rating.
\textbf{Categories:} A Popular/Narrative \quad B Undergraduate Textbooks \quad
C Graduate Textbooks \quad D Monographs \quad E Primary Papers \quad
F Reviews \quad G Historical Collections
\end{tcolorbox}

%--- A: Popular -----------------------------------------------
\subsection*{A \quad Popular Science \& Divulgation}
\addcontentsline{toc}{subsection}{A: Popular Science}
\begin{tcolorbox}[catbox]
\begin{enumerate}[leftmargin=2.5em,label=\textbf{[A\arabic*]},itemsep=3pt]
  \item \tierbadge{t1green}{T1} Al-Khalili, J. \textit{Quantum: A Guide for the Perplexed.} Weidenfeld \& Nicolson, 2003. [N] \recommended
  \item \tierbadge{t1green}{T1} Cox, B.\ \& Forshaw, J. \textit{The Quantum Universe.} Da Capo Press, 2011. [N] \recommended
  \item \tierbadge{t1green}{T1} Feynman, R.~P. \textit{QED: The Strange Theory of Light and Matter.} Princeton University Press, 1985. [N] \essential
  \item \tierbadge{t1green}{T1} Feynman, R.~P., Leighton, R.~B.\ \& Sands, M. \textit{The Feynman Lectures on Physics, Vol.~III.} Addison-Wesley, 1965. [N/T] \essential
  \item \tierbadge{t1green}{T1} Gamow, G. \textit{Thirty Years That Shook Physics.} Dover, 1966. [N] \recommended
  \item \tierbadge{t1green}{T1} Hey, T.\ \& Walters, P. \textit{The New Quantum Universe.} Cambridge University Press, 2003. [N] \useful
  \item \tierbadge{t1green}{T1} Kumar, M. \textit{Quantum: Einstein, Bohr, and the Great Debate.} Norton, 2008. [N] \recommended
  \item \tierbadge{t1green}{T1} Penrose, R. \textit{The Emperor's New Mind.} Oxford University Press, 1989. [N] \useful
  \item \tierbadge{t1green}{T1} Rae, A.~I.~M. \textit{Quantum Physics: Illusion or Reality?} Cambridge University Press, 2004. [N] \recommended
  \item \tierbadge{t1green}{T1} Rovelli, C. \textit{Helgoland.} Riverhead Books, 2021. [N] \useful
\end{enumerate}
\end{tcolorbox}

%--- B: Undergraduate Textbooks --------------------------------
\subsection*{B \quad Undergraduate Textbooks (Tiers 2--3)}
\addcontentsline{toc}{subsection}{B: Undergraduate Textbooks}
\begin{tcolorbox}[catbox]
\begin{enumerate}[leftmargin=2.5em,label=\textbf{[B\arabic*]},itemsep=3pt]
  \item \tierbadge{t2blue}{T2} Atkins, P.~W.\ \& Friedman, R. \textit{Molecular Quantum Mechanics} (5th ed.). Oxford University Press, 2011. [T]
  \item \tierbadge{t2blue}{T2} Binney, J.\ \& Skinner, D. \textit{The Physics of Quantum Mechanics.} Oxford University Press, 2014. [T] \recommended
  \item \tierbadge{t2blue}{T2} Boas, M.~L. \textit{Mathematical Methods in the Physical Sciences} (3rd ed.). Wiley, 2006. [T] \recommended
  \item \tierbadge{t2blue}{T2} Byron, F.~W.\ \& Fuller, R.~W. \textit{Mathematics of Classical and Quantum Physics.} Dover, 1992. [T] \recommended
  \item \tierbadge{t3purple}{T3} Dirac, P.~A.~M. \textit{The Principles of Quantum Mechanics} (4th ed.). Oxford University Press, 1958. [T/P] \essential
  \item \tierbadge{t2blue}{T2} French, A.~P.\ \& Taylor, E.~F. \textit{An Introduction to Quantum Physics.} Norton, 1978. [T] \recommended
  \item \tierbadge{t3purple}{T3} Gasiorowicz, S. \textit{Quantum Physics} (3rd ed.). Wiley, 2003. [T] \recommended
  \item \tierbadge{t2blue}{T2} Griffiths, D.~J.\ \& Schroeter, D.~F. \textit{Introduction to Quantum Mechanics} (3rd ed.). Cambridge University Press, 2018. [T] \essential
  \item \tierbadge{t3purple}{T3} Hannabuss, K. \textit{An Introduction to Quantum Theory.} Oxford University Press, 1997. [T] \recommended
  \item \tierbadge{t3purple}{T3} Isham, C.~J. \textit{Lectures on Quantum Theory.} Imperial College Press, 1995. [T] \recommended
  \item \tierbadge{t2blue}{T2} Jensen, F. \textit{Introduction to Computational Chemistry} (3rd ed.). Wiley, 2017. [T] \recommended
  \item \tierbadge{t2blue}{T2} Levine, I.~N. \textit{Quantum Chemistry} (7th ed.). Pearson, 2014. [T] \recommended
  \item \tierbadge{t2blue}{T2} McMahon, D. \textit{Quantum Mechanics Demystified} (2nd ed.). McGraw-Hill, 2005. [T] \useful
  \item \tierbadge{t3purple}{T3} Messiah, A. \textit{Quantum Mechanics} (2 vols.). North-Holland, 1961--1962 (Dover, 2014). [T] \essential
  \item \tierbadge{t3purple}{T3} Peres, A. \textit{Quantum Theory: Concepts and Methods.} Kluwer, 1993. [T] \recommended
  \item \tierbadge{t2blue}{T2} Phillips, A.~C. \textit{Introduction to Quantum Mechanics.} Wiley, 2003. [T] \useful
  \item \tierbadge{t3purple}{T3} Auletta, G., Fortunato, M.\ \& Parisi, G. \textit{Quantum Mechanics.} Cambridge University Press, 2009. [T] \recommended
  \item \tierbadge{t3purple}{T3} Sakurai, J.~J.\ \& Napolitano, J. \textit{Modern Quantum Mechanics} (3rd ed.). Cambridge University Press, 2020. [T] \essential
  \item \tierbadge{t2blue}{T2} Shankar, R. \textit{Principles of Quantum Mechanics} (2nd ed.). Springer, 1994. [T] \essential
  \item \tierbadge{t2blue}{T2} Strang, G. \textit{Introduction to Linear Algebra} (5th ed.). Wellesley--Cambridge Press, 2016. [T] \essential
  \item \tierbadge{t2blue}{T2} Townsend, J.~S. \textit{A Modern Approach to Quantum Mechanics} (2nd ed.). University Science Books, 2012. [T] \recommended
\end{enumerate}
\end{tcolorbox}

%--- C: Graduate Textbooks ------------------------------------
\subsection*{C \quad Graduate Textbooks (Tiers 4--5)}
\addcontentsline{toc}{subsection}{C: Graduate Textbooks}
\begin{tcolorbox}[catbox]
\begin{enumerate}[leftmargin=2.5em,label=\textbf{[C\arabic*]},itemsep=3pt]
  \item \tierbadge{t5red}{T5} Abrikosov, A.~A., Gorkov, L.~P.\ \& Dzyaloshinski, I.~E. \textit{Methods of Quantum Field Theory in Statistical Physics.} Dover, 1975. [M] \recommended
  \item \tierbadge{t5red}{T5} Altland, A.\ \& Simons, B.~D. \textit{Condensed Matter Field Theory} (2nd ed.). Cambridge University Press, 2010. [T] \essential
  \item \tierbadge{t4orange}{T4} Ballentine, L.~E. \textit{Quantum Mechanics: A Modern Development} (2nd ed.). World Scientific, 2014. [T] \essential
  \item \tierbadge{t5red}{T5} Blank, J., Exner, P.\ \& Havl\'{\i}\v{c}ek, M. \textit{Hilbert Space Operators in Quantum Physics} (2nd ed.). Springer, 2008. [M] \essential
  \item \tierbadge{t4orange}{T4} Bruus, H.\ \& Flensberg, K. \textit{Introduction to Many-Body Quantum Theory.} Oxford University Press, 2004. [T] \essential
  \item \tierbadge{t4orange}{T4} Cohen-Tannoudji, C., Diu, B.\ \& Lalo\"e, F. \textit{Quantum Mechanics} (2 vols., new ed.). Wiley, 2019. [T] \recommended
  \item \tierbadge{t4orange}{T4} Feynman, R.~P.\ \& Hibbs, A.~R. \textit{Quantum Mechanics and Path Integrals} (emended ed.). Dover, 2010. [T] \essential
  \item \tierbadge{t4orange}{T4} Fetter, A.~L.\ \& Walecka, J.~D. \textit{Quantum Theory of Many-Particle Systems.} Dover, 2003. [T] \essential
  \item \tierbadge{t4orange}{T4} Galindo, A.\ \& Pascual, P. \textit{Quantum Mechanics} (2 vols.). Springer, 1990--1991. [T] \essential
  \item \tierbadge{t4orange}{T4} Georgi, H. \textit{Lie Algebras in Particle Physics} (2nd ed.). Westview Press, 1999. [T] \recommended
  \item \tierbadge{t4orange}{T4} Helgaker, T., J\o rgensen, P.\ \& Olsen, J. \textit{Molecular Electronic-Structure Theory.} Wiley, 2000. [M] \essential
  \item \tierbadge{t4orange}{T4} Huang, K. \textit{Statistical Mechanics} (2nd ed.). Wiley, 1987. [T] \recommended
  \item \tierbadge{t4orange}{T4} Koch, W.\ \& Holthausen, M.~C. \textit{A Chemist's Guide to Density Functional Theory} (2nd ed.). Wiley-VCH, 2001. [T] \recommended
  \item \tierbadge{t4orange}{T4} Le Bellac, M. \textit{Quantum Physics.} Cambridge University Press, 2006. [T] \recommended
  \item \tierbadge{t4orange}{T4} Mahan, G.~D. \textit{Many-Particle Physics} (3rd ed.). Kluwer, 2000. [T] \recommended
  \item \tierbadge{t4orange}{T4} Martin, R.~M. \textit{Electronic Structure.} Cambridge University Press, 2004. [T] \recommended
  \item \tierbadge{t4orange}{T4} Martin, R.~M., Reining, L.\ \& Ceperley, D. \textit{Interacting Electrons.} Cambridge University Press, 2016. [M] \essential
  \item \tierbadge{t5red}{T5} Peskin, M.~E.\ \& Schroeder, D.~V. \textit{An Introduction to Quantum Field Theory.} Addison-Wesley, 1995. [T] \recommended
  \item \tierbadge{t5red}{T5} Rudin, W. \textit{Functional Analysis} (2nd ed.). McGraw-Hill, 1991. [M] \essential
  \item \tierbadge{t4orange}{T4} Sachdev, S. \textit{Quantum Phase Transitions} (2nd ed.). Cambridge University Press, 2011. [T] \recommended
  \item \tierbadge{t4orange}{T4} Schwabl, F. \textit{Advanced Quantum Mechanics} (4th ed.). Springer, 2008. [T] \recommended
  \item \tierbadge{t4orange}{T4} Stefanucci, G.\ \& van Leeuwen, R. \textit{Nonequilibrium Many-Body Theory of Quantum Systems.} Cambridge University Press, 2013. [T] \essential
  \item \tierbadge{t4orange}{T4} Szabo, A.\ \& Ostlund, N.~S. \textit{Modern Quantum Chemistry.} Dover, 1996. [T] \essential
  \item \tierbadge{t4orange}{T4} Weinberg, S. \textit{Lectures on Quantum Mechanics} (2nd ed.). Cambridge University Press, 2015. [T] \essential
  \item \tierbadge{t4orange}{T4} Wigner, E.~P. \textit{Group Theory and Its Application to the Quantum Mechanics of Atomic Spectra.} Academic Press, 1959. [T/M] \essential
  \item \tierbadge{t4orange}{T4} Zinn-Justin, J. \textit{Quantum Field Theory and Critical Phenomena} (4th ed.). Oxford University Press, 2002. [M] \recommended
\end{enumerate}
\end{tcolorbox}

%--- D: Research Monographs -----------------------------------
\subsection*{D \quad Research Monographs (Tier 5)}
\addcontentsline{toc}{subsection}{D: Research Monographs}
\begin{tcolorbox}[catbox]
\begin{enumerate}[leftmargin=2.5em,label=\textbf{[D\arabic*]},itemsep=3pt]
  \item \tierbadge{t4orange}{T4}\tierbadge{t5red}{T5} Biedenharn, L.~C.\ \& Louck, J.~D. \textit{Angular Momentum in Quantum Physics} (2 vols.). Addison-Wesley, 1981. [M] \recommended
  \item \tierbadge{t5red}{T5} Breuer, H.-P.\ \& Petruccione, F. \textit{The Theory of Open Quantum Systems.} Oxford University Press, 2002. [M] \essential
  \item \tierbadge{t5red}{T5} Busch, P., Lahti, P.\ \& Mittelstaedt, P. \textit{The Quantum Theory of Measurement} (2nd ed.). Springer, 1996. [M] \recommended
  \item \tierbadge{t5red}{T5} Carmichael, H.~J. \textit{An Open Systems Approach to Quantum Optics.} Springer, 1993. [M] \recommended
  \item \tierbadge{t5red}{T5} Dreizler, R.~M.\ \& Gross, E.~K.~U. \textit{Density Functional Theory.} Springer, 1990. [M] \recommended
  \item \tierbadge{t5red}{T5} Gel'fand, I.~M.\ \& Vilenkin, N.~Ya. \textit{Generalized Functions, Vol.~4.} Academic Press, 1964. [M] \recommended
  \item \tierbadge{t5red}{T5} Haag, R. \textit{Local Quantum Physics} (2nd ed.). Springer, 1996. [M] \recommended
  \item \tierbadge{t5red}{T5} Kadanoff, L.~P.\ \& Baym, G. \textit{Quantum Statistical Mechanics.} Perseus, 1989. [M] \essential
  \item \tierbadge{t5red}{T5} Kraus, K. \textit{States, Effects, and Operations.} Springer, 1983. [M] \recommended
  \item \tierbadge{t5red}{T5} Leggett, A.~J. \textit{Quantum Liquids.} Oxford University Press, 2006. [M] \essential
  \item \tierbadge{t5red}{T5} Lieb, E.~H.\ \& Loss, M. \textit{Analysis} (2nd ed.). AMS, 2001. [M] \recommended
  \item \tierbadge{t5red}{T5} Negele, J.~W.\ \& Orland, H. \textit{Quantum Many-Particle Systems.} Perseus, 1998. [M] \essential
  \item \tierbadge{t5red}{T5} Nielsen, M.~A.\ \& Chuang, I.~L. \textit{Quantum Computation and Quantum Information.} Cambridge University Press, 2000. [M] \essential
  \item \tierbadge{t5red}{T5} Parr, R.~G.\ \& Yang, W. \textit{Density-Functional Theory of Atoms and Molecules.} Oxford University Press, 1989. [M] \essential
  \item \tierbadge{t5red}{T5} Reed, M.\ \& Simon, B. \textit{Methods of Modern Mathematical Physics} (4 vols.). Academic Press, 1972--1979. [M] \essential
  \item \tierbadge{t5red}{T5} Thirring, W. \textit{A Course in Mathematical Physics} (4 vols.). Springer, 1978--1983. [M] \recommended
  \item \tierbadge{t3purple}{T3}\tierbadge{t5red}{T5} von Neumann, J. \textit{Mathematical Foundations of Quantum Mechanics.} Princeton University Press, 1955. [M] \essential
\end{enumerate}
\end{tcolorbox}

%--- E: Primary Papers ----------------------------------------
\subsection*{E \quad Primary Research Papers}
\addcontentsline{toc}{subsection}{E: Primary Research Papers}
\begin{tcolorbox}[catbox]
\begin{enumerate}[leftmargin=2.5em,label=\textbf{[E\arabic*]},itemsep=3pt]
  \item \tierbadge{t5red}{T5} Aharonov, Y.\ \& Bohm, D. Significance of electromagnetic potentials in the quantum theory. \textit{Phys.~Rev.} \textbf{115}, 485 (1959). [P] \essential
  \item \tierbadge{t5red}{T5} Aspect, A., Grangier, P.\ \& Roger, G. Experimental tests of Bell's inequalities using time-varying analyzers. \textit{Phys.~Rev.~Lett.} \textbf{49}, 1804 (1982). [P] \essential
  \item \tierbadge{t5red}{T5} Bardeen, J., Cooper, L.~N.\ \& Schrieffer, J.~R. Theory of superconductivity. \textit{Phys.~Rev.} \textbf{108}, 1175 (1957). [P] \essential
  \item \tierbadge{t4orange}{T4} Becke, A.~D. Density-functional thermochemistry. III. \textit{J.~Chem.~Phys.} \textbf{98}, 5648 (1993). [P] \recommended
  \item \tierbadge{t3purple}{T3} Bell, J.~S. On the Einstein Podolsky Rosen paradox. \textit{Physics} \textbf{1}, 195 (1964). [P] \essential
  \item \tierbadge{t4orange}{T4} Berry, M.~V. Quantal phase factors accompanying adiabatic changes. \textit{Proc.~R.~Soc.~Lond.~A} \textbf{392}, 45 (1984). [P] \essential
  \item \tierbadge{t3purple}{T3} Born, M. Zur Quantenmechanik der Sto\ss vorg\"ange. \textit{Z.~Phys.} \textbf{37}, 863 (1926). [P] \essential
  \item \tierbadge{t3purple}{T3} Born, M.\ \& Jordan, P. Zur Quantenmechanik. \textit{Z.~Phys.} \textbf{34}, 858 (1925). [P] \recommended
  \item \tierbadge{t5red}{T5} Dirac, P.~A.~M. The quantum theory of the emission and absorption of radiation. \textit{Proc.~R.~Soc.~Lond.~A} \textbf{114}, 243 (1927). [P] \recommended
  \item \tierbadge{t3purple}{T3} Einstein, A., Podolsky, B.\ \& Rosen, N. Can quantum-mechanical description of physical reality be considered complete? \textit{Phys.~Rev.} \textbf{47}, 777 (1935). [P] \recommended
  \item \tierbadge{t5red}{T5} Gleason, A.~M. Measures on the closed subspaces of a Hilbert space. \textit{J.~Math.~Mech.} \textbf{6}, 885 (1957). [P] \essential
  \item \tierbadge{t5red}{T5} Gorini, V., Kossakowski, A.\ \& Sudarshan, E.~C.~G. Completely positive dynamical semigroups of N-level systems. \textit{J.~Math.~Phys.} \textbf{17}, 821 (1976). [P] \essential
  \item \tierbadge{t3purple}{T3} Heisenberg, W. \"Uber quantentheoretische Umdeutung. \textit{Z.~Phys.} \textbf{33}, 879 (1925). [P] \recommended
  \item \tierbadge{t3purple}{T3} Heisenberg, W. \"Uber den anschaulichen Inhalt. \textit{Z.~Phys.} \textbf{43}, 172 (1927). [P] \recommended
  \item \tierbadge{t5red}{T5} Heisenberg, W. Zur Theorie des Ferromagnetismus. \textit{Z.~Phys.} \textbf{49}, 619 (1928). [P] \recommended
  \item \tierbadge{t4orange}{T4} Hohenberg, P.\ \& Kohn, W. Inhomogeneous electron gas. \textit{Phys.~Rev.} \textbf{136}, B864 (1964). [P] \essential
  \item \tierbadge{t5red}{T5} Keldysh, L.~V. Diagram technique for nonequilibrium processes. \textit{Sov.~Phys.~JETP} \textbf{20}, 1018 (1965). [P] \essential
  \item \tierbadge{t4orange}{T4} Koch, J.\ et al. Charge-insensitive qubit design. \textit{Phys.~Rev.~A} \textbf{76}, 042319 (2007). [R] \recommended
  \item \tierbadge{t4orange}{T4} Kohn, W.\ \& Sham, L.~J. Self-consistent equations. \textit{Phys.~Rev.} \textbf{140}, A1133 (1965). [P] \essential
  \item \tierbadge{t4orange}{T4} Kubo, R. Statistical-mechanical theory of irreversible processes. \textit{J.~Phys.~Soc.~Japan} \textbf{12}, 570 (1957). [P] \essential
  \item \tierbadge{t5red}{T5} Lindblad, G. On the generators of quantum dynamical semigroups. \textit{Commun.~Math.~Phys.} \textbf{48}, 119 (1976). [P] \essential
  \item \tierbadge{t5red}{T5} Meir, Y.\ \& Wingreen, N.~S. Landauer formula for the current. \textit{Phys.~Rev.~Lett.} \textbf{68}, 2512 (1992). [P] \recommended
  \item \tierbadge{t4orange}{T4} Perdew, J.~P., Burke, K.\ \& Ernzerhof, M. Generalized gradient approximation made simple. \textit{Phys.~Rev.~Lett.} \textbf{77}, 3865 (1996). [P] \recommended
  \item \tierbadge{t4orange}{T4} Perdew, J.~P.\ \& Zunger, A. Self-interaction correction. \textit{Phys.~Rev.~B} \textbf{23}, 5048 (1981). [P] \recommended
  \item \tierbadge{t4orange}{T4} Pulay, P. Ab initio calculation of force constants. \textit{Mol.~Phys.} \textbf{17}, 197 (1969). [P] \recommended
  \item \tierbadge{t4orange}{T4} Runge, E.\ \& Gross, E.~K.~U. Density-functional theory for time-dependent systems. \textit{Phys.~Rev.~Lett.} \textbf{52}, 997 (1984). [P] \recommended
  \item \tierbadge{t3purple}{T3} Schr\"odinger, E. Quantisierung als Eigenwertproblem (Parts I--IV). \textit{Ann.~Phys.} (1926). [P] \recommended
  \item \tierbadge{t5red}{T5} Stone, M.~H. On one-parameter unitary groups in Hilbert space. \textit{Ann.~Math.} \textbf{33}, 643 (1932). [P] \essential
  \item \tierbadge{t3purple}{T3} Wigner, E.~P. On unitary representations of the inhomogeneous Lorentz group. \textit{Ann.~Math.} \textbf{40}, 149 (1939). [P] \useful
  \item \tierbadge{t5red}{T5} Wigner, E.~P. The unreasonable effectiveness of mathematics. \textit{Commun.~Pure Appl.~Math.} \textbf{13}, 1 (1960). [P] \useful
  \item \tierbadge{t5red}{T5} Zurek, W.~H. Pointer basis of quantum apparatus. \textit{Phys.~Rev.~D} \textbf{24}, 1516 (1981). [P] \recommended
\end{enumerate}
\end{tcolorbox}

%--- F: Reviews -----------------------------------------------
\subsection*{F \quad Review Articles \& Surveys}
\addcontentsline{toc}{subsection}{F: Review Articles}
\begin{tcolorbox}[catbox]
\begin{enumerate}[leftmargin=2.5em,label=\textbf{[F\arabic*]},itemsep=3pt]
  \item \tierbadge{t5red}{T5} Rotter, S.\ \& Fan, S. Optical cavities during the age of quantum information. \textit{Science} \textbf{358}, 1444 (2017). [R] \useful
  \item \tierbadge{t5red}{T5} Zurek, W.~H. Decoherence, einselection, and the quantum origins of the classical. \textit{Rev.~Mod.~Phys.} \textbf{75}, 715 (2003). [R] \essential
\end{enumerate}
\end{tcolorbox}

%--- G: Historical Collections --------------------------------
\subsection*{G \quad Historical Collections \& Source Books}
\addcontentsline{toc}{subsection}{G: Historical Collections}
\begin{tcolorbox}[catbox]
\begin{enumerate}[leftmargin=2.5em,label=\textbf{[G\arabic*]},itemsep=3pt]
  \item \tierbadge{t3purple}{T3} Bell, J.~S. \textit{Speakable and Unspeakable in Quantum Mechanics} (2nd ed.). Cambridge University Press, 2004. [M] \essential
  \item \tierbadge{t3purple}{T3} van der Waerden, B.~L.\ (ed.). \textit{Sources of Quantum Mechanics.} Dover, 1967. [M] \recommended
  \item \tierbadge{t3purple}{T3} Wheeler, J.~A.\ \& Zurek, W.~H.\ (eds.). \textit{Quantum Theory and Measurement.} Princeton University Press, 1983. [M] \recommended
  \item \tierbadge{t4orange}{T4} Jammer, M. \textit{The Conceptual Development of Quantum Mechanics.} McGraw-Hill, 1966. [M] \recommended
  \item \tierbadge{t4orange}{T4} Mehra, J.\ \& Rechenberg, H. \textit{The Historical Development of Quantum Theory} (6 vols.). Springer, 1982--2001. [M] \useful
\end{enumerate}
\end{tcolorbox}

\bigskip
\subsection*{Tier--Category Quick-Index}
\addcontentsline{toc}{subsection}{Tier--Category Quick-Index}
\begin{center}
\begin{tabular}{lp{11.5cm}}
\toprule
\textbf{Tier} & \textbf{Recommended categories and key items}\\
\midrule
\tierbadge{t1green}{T1} HS
  & Cat.~A (all); Cat.~B: [B8, B19]; no primary papers\\[3pt]
\tierbadge{t2blue}{T2} BegUG
  & Cat.~B: [B1--B4, B6, B8, B9, B11--B13, B19--B21]; Cat.~A as context;
    Cat.~E: [E7] abstract only\\[3pt]
\tierbadge{t3purple}{T3} AdvUG
  & Cat.~B (all); Cat.~G: [G1--G3];
    Cat.~E: [E5, E7, E8, E10, E13, E14, E27] in full\\[3pt]
\tierbadge{t4orange}{T4} MSc
  & Cat.~C: [C3, C5, C7--C11, C13, C16, C22--C26]; Cat.~D: [D8, D12, D14];
    Cat.~E: [E1, E6, E7, E11, E16, E19, E20, E23] as primary reading;
    Cat.~G: [G1, G3, G4]\\[3pt]
\tierbadge{t5red}{T5} PhD
  & Cat.~D (all); Cat.~E (complete); Cat.~F (all); Cat.~C as reference;
    current arXiv preprints\\[3pt]
Research
  & Cat.~D (all); Cat.~E (all); Cat.~F (all); current arXiv; journal club\\
\bottomrule
\end{tabular}
\end{center}

%=============================================================
\end{document}
